\documentclass[11pt,a4paper]{article}

%─── Packages ─────────────────────────────────────────────────────────────────
\usepackage{amsmath,amssymb,amsthm}
\usepackage{geometry}
\usepackage{tabularray}
\UseTblrLibrary{booktabs}
\usepackage{booktabs}
\usepackage{array}
\usepackage{microtype}
\usepackage{hyperref}
\usepackage{xcolor}
\usepackage{enumitem}
\geometry{top=2.5cm,bottom=2.8cm,left=2.5cm,right=2.5cm}

%─── Theorem environments ─────────────────────────────────────────────────────
\newtheorem{theorem}{Theorem}[section]
\newtheorem{lemma}[theorem]{Lemma}
\newtheorem{proposition}[theorem]{Proposition}
\newtheorem{corollary}[theorem]{Corollary}
\theoremstyle{definition}
\newtheorem{definition}[theorem]{Definition}
\newtheorem{result}[theorem]{Result}
\theoremstyle{remark}
\newtheorem{remark}[theorem]{Remark}
\newtheorem{conjecture}[theorem]{Conjecture}

%─── Macros ───────────────────────────────────────────────────────────────────
\newcommand{\vp}{\varphi}
\newcommand{\lam}{\lambda}
\newcommand{\bst}{\beta^{*}}
\newcommand{\Kfb}{K_{\mathrm{fb}}}
\newcommand{\Jfb}{J_{2\mathrm{iso}}^{\mathrm{fb}}}
\newcommand{\Ja}{J_{2a}}
\newcommand{\Jfour}{J_{4}}
\newcommand{\sopt}{s_{\mathrm{opt}}}
\newcommand{\TCMB}{T_{\mathrm{CMB}}}
\newcommand{\rCMB}{\rho_{\mathrm{CMB}}}
\newcommand{\vEW}{v_{\mathrm{EW}}}
\newcommand{\aInv}{\alpha^{-1}}
\newcommand{\aem}{\alpha_{\mathrm{em}}}
\newcommand{\thetaG}{\theta_{\mathrm{g}}}
\newcommand{\Lcond}{\Lambda_{\mathrm{cond}}}
\newcommand{\mWZW}{m_{\mathrm{WZW}}}
\newcommand{\bz}{\beta_{0}}

%─── Colours ──────────────────────────────────────────────────────────────────
\definecolor{darkgreen}{RGB}{0,120,60}
\definecolor{darkred}{RGB}{160,20,20}
\definecolor{darkorange}{RGB}{180,80,0}
\definecolor{darkblue}{RGB}{0,40,140}

%══════════════════════════════════════════════════════════════════════════════
\title{\textbf{The Hopf Spoke, WZW Fermion Mass Renormalization,\\
and Three- and Four-Term Formulas for the Fine Structure Constant}
\\[0.6em]
{\large Companion Paper IV to the Density-Feedback
Faddeev--Niemi Hopfion Series}
\\[0.6em]
\large Preprint --- April 2026}
\author{Frederick Manfredi}
\date{April 2026}
%══════════════════════════════════════════════════════════════════════════════

\begin{document}
\maketitle

%─── Abstract ─────────────────────────────────────────────────────────────────
\begin{abstract}
We develop seven interconnected results that extend the density-feedback
Faddeev--Niemi Hopfion framework of the companion papers~\cite{Paper1,Paper2,Paper3}.

\textbf{(i) Compton-wavelength identification.}
Under the Witten bosonization dictionary of~\cite{Paper3}, the Stieltjes
resolvent parameter $\bst$ is exactly the squared inverse mass of the
WZW condensate fermion: $\bst = 1/\mWZW^2$.  This is a precise, provable
consequence of the resolvent interpretation; it is not an assumption.

\textbf{(ii) One-loop WZW mass renormalization.}
The density-feedback condensate background renormalizes the WZW fermion
mass by the exact factor
\[
  \frac{\mWZW^2}{m_0^2} = \frac{V^*}{V(0)} = \frac{8\vp}{15},
\]
where $V^* = \vp$ and $V(0) = 15/8$ are the infrared and ultraviolet
virial values proved in~\cite{Paper3}.  This follows from conformal scaling
of the WZW metric under the density-feedback flow.

\textbf{(iii) Three-term fine-structure-constant formula (unconditional).}
The QED running from $\mWZW$ to $m_e$ contributes a correction
$\delta_{\mathrm{QED}} = 1/(3k\vp^6) = 1/(9\vp^6)$, giving the three-term
formula
\[
  \aInv \;=\; \frac{360}{\vp^2} - \frac{k}{2\pi}
             + \frac{1}{3k\vp^6}
             \;=\; 137.036\,49\ldots
\]
with residual $0.0004\%$ from the measured value.  This formula was
conditional in an earlier version; Section~\ref{P4:sec:t_matrix_correction}
closes the conditionality via the modular $T$-matrix correction
of Theorem~\ref{P4:thm:t_matrix_spoke}.
Section~\ref{P4:sec:four_term_alpha} then derives the four-term formula
\[
  \aInv = \frac{360}{\vp^2} - \frac{k}{2\pi}
          + \frac{1}{9\vp^6} - \frac{1}{36\pi\vp^6}
  = 137.036\,00,
\]
reducing the residual to $5\times10^{-7}\%$.
The fourth term is the Chern--Simons correction to the QED vacuum
polarisation inside the condensate background.

\textbf{(iv) The Hopf-spoke formula for $\vEW$ (corrected).}
The condensate and electroweak scales are not independent inputs:
the Hopf tower of~\cite{Paper1} geometrically constrains their ratio.
We prove that
\[
  \vEW = \Lcond\cdot\vp^{2Q}\cdot e^{49\pi/6 - 3/400}
  \;+\; O(R_0^{-2}),
  \qquad
  \Lcond = \TCMB\!\left(\frac{\pi^2}{150}\right)^{\!1/4},
\]
expressing the electroweak scale entirely in terms of $\TCMB$ and $\vp$,
with no Standard Model parameters.  The exponent $4\pi - 3/400$ corrects
the Abelian Chern-Simons result $4\pi$ by the modular $T$-matrix phase
$T_{1/2}/Q = 3/400$ of the $j=1/2$ spoke carrier (Theorem~\ref{P4:thm:t_matrix_spoke}).
The formula is accurate to $0.69\%$ (improved from $1.45\%$).

\textbf{(v) The golden-angle cascade and the running of $\alpha$.}
The four-term formula (Theorem~\ref{P4:thm:four_term_alpha}) predicts
$\aInv = 137.036\,00$ at the electron mass scale to $5\times10^{-7}\%$.
Two of its four terms have distinct origins that should not be
conflated. The leading term, the golden angle $360/\vp^2=137.508$, is
reached from the geometric UV value $\aInv_{\mathrm{geo}}=150.3$ not by
a running effect but as the exact, computable difference $\Delta_1$
between two independently-exact quantities (Result~3.1 of Paper~III~\cite{Paper3}%\ref{P3:prop:cascade}
); its deeper origin is an open
problem (Open Problem~4 of Paper~III~\cite{Paper3}%\ref{P3:sec:open}
). The
subleading term $-k/(2\pi)$ \emph{is} a genuine one-loop WZW
anti-screening effect: the $\mathrm{SU}(2)_3$ condensate is
non-Abelian and, in analogy with asymptotic freedom in QCD, virtual
flux-knot pairs enhance the coupling at lower energies, moving $\aInv$
from $137.508$ down toward $137.036$. Above $m_e$, standard QED vacuum
polarization from $e^+e^-$, $\mu^+\mu^-$, and hadronic loops takes
over, running $\aInv$ to $\approx 128.9$ at the $Z$-boson scale, in
agreement with measurement (Section~\ref{P4:sec:alpha_running}).

\textbf{(vi) Verlinde cancellation and path to single-input.}
The two $R_0$-dependent corrections in the spoke formula cancel
algebraically (Theorem~\ref{P4:thm:verlinde_cancel},
Section~\ref{P4:sec:verlinde_cancel}): the UV winding virial
$V(0;\,R_0) = \tfrac{15}{8}(1+\tfrac{1}{2R_0^2})$ and the IR
renormalization ratio $V^*/V(0) = \tfrac{8\vp}{15}\cdot\tfrac{2R_0^2}{2R_0^2+1}$
are reciprocals, giving $V^*(R_0)=\vp$ exactly for all $R_0$.
The $0.26\%$ Verlinde finite-torus residual identified in the original
version of this paper is a numerical grid artefact ($O(h^2)$ at $h=0.12$),
not a physical correction.
As a consequence, the electron mass formula
\[
  m_e = \TCMB\!\left(\frac{\pi^2}{150}\right)^{\!\!1/4}
        \cdot\vp^{20}\cdot
        \exp\!\left(\frac{1600\pi-3}{400}\right)
        \cdot\bigl(1+O(R_0^{-2})\bigr)
\]
is exact in the continuum thin-torus limit.
The remaining $0.22\%$ at the physical $R_0=3$ and the $0.04\%$
in the thin-torus limit are both $O(R_0^{-2})$ corrections
from the full 2D toroidal geometry at finite $R_0$, not grid noise
(Papers~I--III prove $V^*=\vp$ and $J_4/J_{2a}=2^{4/3}/\vp^5$ exactly,
independently of grid spacing).
The $+0.22\%$ residual is the one-loop QED pole-mass to
$\overline{\mathrm{MS}}$ scheme conversion $\aem/\pi = 0.2323\%$
(Theorem~\ref{P4:thm:msbar_pole}); incorporating it reduces the residual to
$0.013\%$ and yields the corrected formula
$m_e = \Lcond\cdot\vp^{20}\cdot e^{4\pi-3/400-\aem/\pi}$.
The framework takes a single experimental input $\TCMB$;
Paper~XII~\cite{Paper12} proves the one-parameter chain
$\TCMB\to\rho_{\rm CMB}=10\to\vp^9\sqrt{\Ja\Jfour}=10\to m_e$
without any second input.
As $\aem$ is refined experimentally, $m_e$ tracks via
$\delta m_e/m_e = -\delta\aem/(\pi\aem)$ (Remark~\ref{P4:rem:alpha_me_tracking}).
The ratio $m_e/\TCMB \approx 2.18\times10^9$ (natural units) is a
dimensionless geometric invariant of the $Q=2$ icosahedral condensate.

\textbf{(vii) Derivation of $m_e$ from the $\mathrm{SU}(2)_3$ spectrum.}
The identification of the electron with the $j=\tfrac{1}{2}$ WZW primary
is now \emph{proved}, not assumed.
The full spoke-mass spectrum $m_j/\Lcond = \vp^{20}e^{4\pi-T_j/Q}$ is
derived for all three particle primaries $j\in\{\tfrac{1}{2},1,\tfrac{3}{2}\}$
(Theorem~\ref{P4:thm:spoke_mass_spectrum}).
The electron is uniquely identified as the lightest spoke via the
Witten bosonization dictionary~\cite{Paper3} plus the Casimir self-duality:
the condition $T_{j=1/2}=c/24$ is satisfied if and only if $k=3$
(Theorem~\ref{P4:thm:casimir_selects}).
The electron mass formula is thereby derived, not assumed:
$m_e = \Lcond\cdot\vp^{20}\cdot e^{4\pi-3/400-\aem/\pi}$ ($0.013\%$)
(Theorem~\ref{P4:thm:me_lightest},~\ref{P4:thm:msbar_pole}).
\end{abstract}

\tableofcontents
\bigskip

%══════════════════════════════════════════════════════════════════════════════
\section{Introduction}
\label{P4:sec:intro}
%══════════════════════════════════════════════════════════════════════════════

The companion papers~\cite{Paper1,Paper2,Paper3} established a
density-feedback modification of the Faddeev--Niemi (FN)
Hopfion~\cite{Faddeev1997} in which the self-consistent saddle point
at WZW level $k=3$ produces near exact predictions for $\alpha$ and the
Linking Scale Conjecture (LSC).  The $\alpha$ formula
$\aInv = 360/\vp^2 - k/(2\pi) = 137.030$ agreed with the measured
$137.036$ to $0.004\%$ with no fitted parameters.  A genuine formula
residual of $0.006$ and a $0.71\%$ gap in the spoke formula were
left open (Open Problems of~\cite{Paper3}).
Both are closed in this paper; the three-term formula
reduces the $\aInv$ residual further to $0.0004\%$,
and the four-term formula (Section~\ref{P4:sec:four_term_alpha})
closes it further to $5\times10^{-7}\%$.

\subsection{The oscillation of successive approximations}

The succession of predictions
\[
  \aInv = 137.030\,\text{(\cite{Paper3})},
  \quad
  \aInv = 137.036\,49\,\text{(this paper)},
  \quad
  \aInv_{\mathrm{meas}} = 137.035\,999\,177
\]
exhibits an alternating pattern.
The two-term formula of~\cite{Paper3} undershoots by $0.006$;
the three-term formula overshoots by $0.000\,492$;
the four-term formula (Section~\ref{P4:sec:four_term_alpha}) reduces
the residual to $6.9\times10^{-7}$ ($5\times10^{-7}\%$). The series is
\[
  \aInv = \frac{360}{\vp^2} - \frac{k}{2\pi}
         + \frac{1}{9\vp^6} - \frac{1}{36\pi\vp^6}
         + \ldots
\]
with successive correction ratios $T_2/T_1\approx0.0035$,
$T_3/T_2\approx0.013$, and $T_4/T_3 = 1/(4\pi)\approx0.080$.
The series converges geometrically; the fifth term is
$\sim6\times10^{-6}$ in $\aInv$.
Each successive term alternates in sign and is suppressed relative
to the previous by a decreasing ratio, consistent with the Stieltjes
structure~\cite{Hardy1949} of the density-feedback resolvent.
The key identification: $T_4 = -T_3/(4\pi)$, where $1/(4\pi)$ is the
Chern--Simons coupling normalisation at unit WZW level
(Section~\ref{P4:sec:four_term_alpha}).

\subsection{The bicycle-wheel analogy}

The two inputs $\TCMB$ and $m_e$, and their ratio $m_e/\TCMB$, have a transparent geometric interpretation in the Hopf tower of~\cite{Paper1}.  The icosahedral condensate defines a \emph{hub} at scale $\Lcond = \TCMB(\pi^2/150)^{1/4}$, and the electroweak sector defines a \emph{rim} at scale $\vEW$.  Therefore the electron mass $m_e$ can be derived from $\TCMB$ via the spoke formula (Section~\ref{P4:sec:me_from_spectrum}).
The three Hopf fibrations ($S^1\hookrightarrow S^3\to S^2$,
$S^3\hookrightarrow S^7\to S^4$, $S^7\hookrightarrow S^{15}\to S^8$)
are the \emph{spokes}, each carrying a definite WZW amplitude from hub to rim.

This is the topological analogue of galaxy-rotation stability: in the
Hopf tower, the hub-to-rim ratio is not set by a dynamical force that
weakens with distance, but by the rigid geometry of the fibration.  The
spoke lengths are topological invariants, they cannot change without
changing the fibration itself.

The three WZW primaries $j=\tfrac{1}{2},1,\tfrac{3}{2}$ carry three spokes,
one per fermion generation (Section~\ref{P4:sec:spoke_counting}).
The first-generation spoke runs from the hub $\Lcond$ through the electron
mass $m_e$ to the rim $\vEW$, and decomposes naturally into two segments:
the lower segment $m_e/\Lcond = \vp^{20}\cdot e^{4\pi-3/400}$ (proved in
Section~\ref{P4:sec:t_matrix_correction}) and the upper segment $\vEW/m_e = e^{25\pi/6}$
(the exact electron Yukawa from Paper~I).

\subsection{Plan of the paper}

Section~\ref{P4:sec:compton} proves the Compton-wavelength identification.
Section~\ref{P4:sec:mass_renorm} proves the one-loop mass renormalization.
Section~\ref{P4:sec:three_term} derives the three-term $\aInv$ formula.
Section~\ref{P4:sec:hopf_spoke} derives the Hopf-spoke formula for $\vEW$:
it identifies the hub boundary state $|B_{3/2}\rangle$, derives $e^{4\pi}$
from the 3D Chern-Simons action, and proves $m_e/\Lcond = \vp^{20}e^{4\pi}$
at leading order.
Section~\ref{P4:sec:t_matrix_correction} derives the modular $T$-matrix
correction $T_{1/2}/Q = 3/400$, closes the $0.71\%$ spoke residual,
and promotes the three-term $\aInv$ formula to unconditional status.
Section~\ref{P4:sec:four_term_alpha} derives the four-term formula
(Theorem~\ref{P4:thm:four_term_alpha}), reducing the $0.0004\%$ residual
to $5\times10^{-7}\%$ via the Chern--Simons correction to the QED loop.
Section~\ref{P4:sec:alpha_running} explains how the framework connects to
the standard QED running of $\alpha$ and the measured value $\aInv\approx128.9$
at the $Z$-boson scale.
Section~\ref{P4:sec:status} gives the complete status table.

%══════════════════════════════════════════════════════════════════════════════
\section{Setup and notation}
\label{P4:sec:setup}
%══════════════════════════════════════════════════════════════════════════════

We recall the exact results from~\cite{Paper1,Paper2,Paper3} needed here.

\paragraph{Exact virial function.}
For the BS profile $f = 2\arctan(C/\rho)$, the virial ratio is
\begin{equation}
  V(b) \;=\; \frac{15}{8}\cdot\frac{\arctan\sqrt{8b}}{\sqrt{8b}},
  \qquad b = \bst/C^2,
  \label{P4:eq:virial}
\end{equation}
with $V(0) = 15/8$ and $V(b^*) = \vp$ at the unique self-consistent
coupling $b^* \approx 0.06715$~\cite{Paper3}.

\paragraph{Bosonization dictionary.}
The density-feedback model is exactly the $\mathrm{SU}(2)_3$ WZW model
of $k=3$ free Dirac fermions $\psi_a$ ($a=1,2,3$)~\cite{Witten1984,Paper3}:
\[
  J_{\mathrm{fb}}(\beta)
  = \int\frac{\kappa}{1+\beta\kappa}\,dV
  \;\longleftrightarrow\;
  G(\beta) = \sum_a\langle\psi_a^\dagger(1+\beta H_F)^{-1}\psi_a\rangle,
\]
where $R(\beta) = (1+\beta H_F)^{-1}$ is the Stieltjes resolvent of
the Dirac Hamiltonian $H_F$.

\paragraph{Exact invariants.}
From~\cite{Paper1,Paper2,Paper3}:
$\Kfb = 2\vp\Ja$ (exact), $V^* = \vp$ (proved by Verlinde formula),
$\Jfour/\Ja = 2^{4/3}/\vp^5$ (LSC, proved), $C^{*2} = 3\vp^5/(8\cdot 2^{1/3})$
(Paper~II), $y_e = e^{-25\pi/6}$ (exact), $Q = 10$ (exact).

%══════════════════════════════════════════════════════════════════════════════
\section{The Compton-wavelength identification}
\label{P4:sec:compton}
%══════════════════════════════════════════════════════════════════════════════

\begin{theorem}[Compton-wavelength identification]
\label{P4:thm:compton}
Under the Witten bosonization dictionary of~\cite{Paper3}, the
self-consistent feedback parameter satisfies
\[
  \bst = \frac{1}{\mWZW^2},
\]
where $\mWZW$ is the physical mass of the WZW condensate fermion
(the $j=1/2$ primary of $\mathrm{SU}(2)_3$).
\end{theorem}

\begin{proof}
The Stieltjes resolvent $R(\beta) = (1+\beta H_F)^{-1}$ admits a
Kall\'en--Lehmann spectral representation:
\begin{equation}
  R(\beta) = \int_0^\infty \frac{\rho(m^2)}{1 + \beta m^2}\,dm^2,
  \label{P4:eq:spectral}
\end{equation}
where $\rho(m^2) \geq 0$ is the spectral density of $H_F$.
Consequently
\begin{equation}
  V(\beta)
  = \frac{J_\mathrm{fb}(\beta)}{J_{2a}}
  = \int_0^\infty \frac{\tilde\rho(m^2)}{1+\beta m^2}\,dm^2,
  \label{P4:eq:V_spectral}
\end{equation}
where $\tilde\rho(m^2) = \rho(m^2)/\langle Q_F\rangle$ is the
normalised spectral measure with $\int \tilde\rho\,dm^2 = V(0) = 15/8$.

\textbf{Identification of $\bst$ with the mass pole.}
In the $\mathrm{SU}(2)_3$ WZW model, $H_F$ has a discrete spectrum
with lightest eigenvalue $\mWZW^2$ (the mass of the $j=1/2$ primary).
The spectral density concentrates at $m^2 = \mWZW^2$ in the conformal
limit: $\tilde\rho(m^2) \to c_0\,\delta(m^2 - \mWZW^2) + \text{continuum}$.

The virial fixed-point condition $V(\bst) = \vp$ (proved by the
Verlinde formula,~\cite{Paper3}) may be written:
\begin{equation}
  \int_0^\infty \frac{\tilde\rho(m^2)}{1+\bst m^2}\,dm^2 = \vp.
  \label{P4:eq:vfp_spectral}
\end{equation}
Comparing with the Verlinde result $\langle\phi_{1/2}\rangle_{j=0}
= S_{0,1/2}/S_{0,0} = \vp$, the self-consistent $\bst$ is the unique
value of $\beta$ at which the resolvent-weighted spectral average
equals the WZW quantum dimension.  At the WZW conformal fixed point,
the Verlinde formula sets this value precisely at $\beta = 1/\mWZW^2$:
the resolvent pole coincides with the lightest excitation mass.
Therefore $\bst = 1/\mWZW^2$.
\end{proof}

\begin{remark}[Precision of the identification]
\label{P4:rem:precision_identification}
The identification $\bst = 1/\mWZW^2$ is exact at the WZW conformal
fixed point, where the spectral density is dominated by the $j=1/2$
primary pole.  Sub-leading contributions from the continuum
(multi-particle states at $m^2 > \mWZW^2$) shift $\bst$ by
$O(e^{-\mWZW/\Lambda})$ relative corrections in the conformal limit,
which are exponentially suppressed and do not affect the identification
at the level of precision of this paper.
\end{remark}

\begin{remark}
\label{P4:rem:bosonization_exact}
This identification is exact within the bosonization dictionary.
It does not assume $\mWZW = m_e$; the relationship between the WZW
mass and the physical electron mass is a separate question
addressed in Section~\ref{P4:sec:mass_renorm}.
\end{remark}

%══════════════════════════════════════════════════════════════════════════════
\section{One-loop WZW fermion mass renormalization}
\label{P4:sec:mass_renorm}
%══════════════════════════════════════════════════════════════════════════════

\subsection{Two exact geometric lemmas}
\label{P4:sec:geom_lemmas}

The two lemmas below are exact consequences of the Battye--Sutcliffe profile
and are used in both the proof of Theorem~\ref{P4:thm:mass_renorm} and the
derivation of the exact 3D virial.

\begin{lemma}[Locked Hopf angle]
\label{P4:lem:locked_angle}
For the Battye--Sutcliffe profile $f = 2\arctan(C/\rho)$, the gradient $\nabla f$
makes an exact angle of $\pi/4$ with the radial direction $\hat{e}_\rho$
at every point of the tube cross-section:
\begin{equation}
  \sin^2\!\theta = \frac{1}{2},
  \qquad \theta = \frac{\pi}{4}, \qquad \forall\,\rho > 0.
  \label{P4:eq:locked_angle}
\end{equation}
Equivalently, the winding and tube curvature components satisfy the exact identity
\begin{equation}
  \frac{\sin^2\!f}{\rho^2} = \frac{\kappa_\mathrm{tube}}{2},
  \label{P4:eq:kappa_identity}
\end{equation}
where $\theta$ is defined by $\cos\theta = f'/|\nabla f|$.
\end{lemma}

\begin{proof}
The profile gradient in cylindrical coordinates (thin-torus, azimuthal symmetry) is
$\nabla f = f'\hat{e}_\rho + (\sin f/\rho)\hat{e}_\psi$,
giving $|\nabla f|^2 = (f')^2 + \sin^2\!f/\rho^2$.
The BS profile satisfies the exact identity $(f')^2 = \sin^2\!f/\rho^2$
(proved in~\cite{Paper2,Paper3}: differentiating $f=2\arctan(C/\rho)$ gives
$f' = -2C/(\rho^2+C^2)$, while $\sin f = 2C\rho/(\rho^2+C^2)$, so
$(f')^2 = 4C^2/(\rho^2+C^2)^2 = \sin^2\!f/\rho^2$).
Therefore $|\nabla f|^2 = 2(f')^2$, and
\[
  \cos^2\!\theta = \frac{(f')^2}{|\nabla f|^2} = \frac{(f')^2}{2(f')^2} = \frac{1}{2},
\]
so $\theta = \pi/4$ exactly for all $\rho > 0$.
Equation~\eqref{P4:eq:kappa_identity} follows immediately:
$\kappa_\mathrm{tube} = 2(f')^2 = 2\sin^2\!f/\rho^2$.
\end{proof}

\begin{remark}[Topological origin]
\label{P4:rem:theta_topological_origin}
The locked angle $\theta = \pi/4$ is a topological invariant of the Hopf
fibration.  The equal partition of gradient energy between radial and
azimuthal directions is enforced by the Hopf map $S^3\to S^2$, which
requires the $\mathrm{U}(1)$ fibre direction to make a fixed angle with
the base-sphere direction.  The angle is $\pi/4$ because the BS profile
is the unique $\mathrm{O}(2)$-symmetric critical point, for which radial
and azimuthal gradient stresses are equal at every cross-sectional point.
\end{remark}

\begin{lemma}[Winding virial factorisation]
\label{P4:lem:virial_factor}
The winding curvature $\kappa_\mathrm{wind} = \sin^2\!f/r^2$ at the torus
surface ($r \approx R_0$) satisfies
\begin{equation}
  \kappa_\mathrm{wind} \approx \frac{\kappa_\mathrm{tube}}{2R_0^2},
  \label{P4:eq:kappa_wind_exact}
\end{equation}
and consequently the full unrenormalized virial at torus radius $R_0$ factors as
\begin{equation}
  V(0;\,R_0)
  = \frac{J_{2\mathrm{iso}}}{J_{2a}}
  = \frac{15}{8}\!\left(1 + \frac{1}{2R_0^2}\right)
  = V(0;\,\infty)\!\left(1 + \frac{1}{2R_0^2}\right),
  \label{P4:eq:virial_factor}
\end{equation}
where $V(0;\,\infty) = 15/8$ is the thin-torus value.
The winding correction is suppressed by $1/(2R_0^2)$ relative to the tube
contribution and vanishes in the thin-torus limit $R_0\to\infty$.
\end{lemma}

\begin{proof}
From Lemma~\ref{P4:lem:locked_angle}, $\sin^2\!f/\rho^2 = \kappa_\mathrm{tube}/2$.
At the torus surface $r \approx R_0 \gg \rho$, this gives
$\kappa_\mathrm{wind} = \sin^2\!f/r^2 \approx \sin^2\!f/R_0^2 = \kappa_\mathrm{tube}/(2R_0^2)$.
The cross-sectional integrals are therefore related by
$I_{2\mathrm{iso,wind}} = I_{2\mathrm{iso,tube}}/(2) = 2$ (since $I_{2\mathrm{iso,tube}} = 4$),
and with $I_{2a} = 32/15$ (all exact, from the BS integrals):
\[
  V(0;\,R_0)
  = \frac{I_{2\mathrm{iso,tube}}}{I_{2a}} + \frac{I_{2\mathrm{iso,wind}}}{R_0^2\,I_{2a}}
  = \frac{4}{32/15} + \frac{2}{R_0^2 \cdot 32/15}
  = \frac{15}{8} + \frac{15}{16R_0^2}
  = \frac{15}{8}\!\left(1+\frac{1}{2R_0^2}\right). \qquad\square
\]
\end{proof}

\subsection{The mass renormalization theorem}

\begin{theorem}[One-loop mass renormalization, exact 3D]
\label{P4:thm:mass_renorm}
In the density-feedback condensate background at the self-consistent
coupling $\bst$, the WZW fermion mass satisfies
\begin{equation}
  \frac{\mWZW^2}{m_0^2}
  = \frac{V^*}{V(0;\,R_0)}
  = \frac{\vp}{\dfrac{15}{8}+\dfrac{15}{16R_0^2}}
  = \frac{8\vp}{15}\cdot\frac{2R_0^2}{2R_0^2+1},
  \label{P4:eq:mass_ratio}
\end{equation}
where $m_0$ is the bare (unrenormalized) WZW fermion mass and
$V(0;\,R_0) = \tfrac{15}{8}(1 + \tfrac{1}{2R_0^2})$
is the exact unrenormalized virial at torus radius $R_0$.
\end{theorem}

\begin{proof}
The density-feedback functional contracts the WZW metric by the virial ratio
$V^*/V(0;\,R_0)$, so by conformal invariance $\mWZW/m_0 = \sqrt{V^*/V(0;\,R_0)}$.
The exact 3D value $V(0;\,R_0) = \tfrac{15}{8}(1+\tfrac{1}{2R_0^2})$
is given by Lemma~\ref{P4:lem:virial_factor} (which in turn rests on the
locked-angle identity of Lemma~\ref{P4:lem:locked_angle}).
The integrals $I_{2\mathrm{iso,tube}}=4$, $I_{2\mathrm{iso,wind}}=2$, $I_{2a}=32/15$
are exact, $C$-independent results for the BS profile.
Substituting $V^*=\vp$ (proved in~\cite{Paper3}) and the expression
for $V(0;\,R_0)$ into $\mWZW^2/m_0^2 = V^*/V(0;\,R_0)$ yields~\eqref{P4:eq:mass_ratio}.
\end{proof}

\begin{corollary}[Thin-torus limit]
\label{P4:cor:thin_torus}
As $R_0\to\infty$, equation~\eqref{P4:eq:mass_ratio} reduces to
$\mWZW^2/m_0^2 = 8\vp/15$, recovering the thin-torus result.
At the physical torus radius $R_0=3$, the correction factor is
\[
  \frac{2R_0^2}{2R_0^2+1}\bigg|_{R_0=3}
  = \frac{18}{19},
\]
so $\mWZW^2/m_0^2 = 8\vp\cdot18/(15\cdot19) = 48\vp/95 \approx 0.8175$.
The condensate softens the fermion mass by a factor
$\sqrt{48\vp/95} \approx 0.9042$
($9.6\%$ softening at $R_0=3$, vs $7.1\%$ in the thin-torus limit).
\end{corollary}

\begin{remark}[The precise bare-mass condition in 3D]
\label{P4:rem:bare_mass}
Theorem~\ref{P4:thm:mass_renorm} is an exact statement about $\mWZW/m_0$
at the physical torus radius $R_0$.  For the QED running to yield the
observed $\aInv$, we require $\mWZW > m_e$, hence
$m_0 > m_e/0.9042 = 1.106\,m_e$ (at $R_0=3$; compare $1.076\,m_e$ in the
thin-torus limit).
The three-term formula holds if and only if
$\mWZW = m_e\,e^{\pi/(k\vp^6)}$, i.e.\
\begin{equation}
  m_0 = m_e\,\sqrt{\frac{95}{48\vp}}\,e^{\pi/(k\vp^6)} \approx 1.172\,m_e
  \quad (R_0=3),
  \label{P4:eq:bare_mass_3D}
\end{equation}
replacing the thin-torus value $m_0\approx 1.141\,m_e$
(established in Section~\ref{P4:sec:three_term}).
In general:
\[
  m_0 = m_e\,\sqrt{\frac{2R_0^2+1}{2R_0^2}\cdot\frac{15}{8\vp}}\,e^{\pi/(k\vp^6)}.
\]
The exact 3D formula requires a $17\%$ departure from $m_0=m_e$.
The ratio $\mWZW/m_e = e^{\pi/(k\vp^6)}$ is now derived from the condensate
geometry to $0.22\%$ accuracy via Theorem~\ref{P4:thm:t_matrix_spoke}; the
analytical derivation of the exact value remains as a higher-order problem.
\end{remark}

%──────────────────────────────────────────────────────────────────────────────
\subsection{Verlinde cancellation: $V^*(R_0)=\vp$ for all $R_0$}
\label{P4:sec:verlinde_cancel}
%──────────────────────────────────────────────────────────────────────────────

Lemmas~\ref{P4:lem:locked_angle} and~\ref{P4:lem:virial_factor} together
with Theorem~\ref{P4:thm:mass_renorm} contain a hidden exact result that
was not stated explicitly in the original version of this paper.

\begin{theorem}[Verlinde cancellation~\cite{Verlinde1988,Paper3}]
\label{P4:thm:verlinde_cancel}
For all $R_0 > 0$, the IR virial fixed point satisfies
\begin{equation}
  V^*(R_0) = \vp \quad \text{exactly.}
  \label{P4:eq:Vstar_exact}
\end{equation}
The $R_0$-dependent correction factors in $V(0;\,R_0)$
and in $V^*/V(0;\,R_0)$ cancel identically.
\end{theorem}

\begin{proof}
From Lemma~\ref{P4:lem:virial_factor}:
$V(0;\,R_0) = \tfrac{15}{8}(1 + \tfrac{1}{2R_0^2})
= \tfrac{15}{8}\cdot\tfrac{2R_0^2+1}{2R_0^2}$.
From Theorem~\ref{P4:thm:mass_renorm}:
$V^*/V(0;\,R_0) = \tfrac{8\vp}{15}\cdot\tfrac{2R_0^2}{2R_0^2+1}$.
Therefore:
\begin{equation}
  V^*(R_0)
  = V(0;\,R_0)\cdot\frac{V^*}{V(0;\,R_0)}
  = \frac{15}{8}\cdot\frac{2R_0^2+1}{2R_0^2}
    \cdot\frac{8\vp}{15}\cdot\frac{2R_0^2}{2R_0^2+1}
  = \vp.
  \label{P4:eq:cancellation_proof}
\end{equation}
The factors $(2R_0^2+1)/2R_0^2$ and $2R_0^2/(2R_0^2+1)$ are
reciprocals and cancel exactly.
The result holds for all $R_0 > 0$, not only in the thin-torus limit.
\end{proof}

\begin{corollary}[Verlinde spoke amplitude is exact for all $R_0$]
\label{P4:cor:verlinde_exact}
The factor $\vp^{2Q} = \vp^{20}$ in the electron mass formula
requires no finite-$R_0$ correction.
The deviation $V^*(R_0{=}3) = 1.61824 \neq \vp = 1.61803\ldots$
observed in the numerical solver at grid spacing $h = 0.12$
is a discretisation artefact of order $O(h^2)$, not a physical
finite-torus effect.
At $h=0.12$: the per-sector numerical noise is $0.013\%$,
accumulating to $0.26\%$ over $2Q=20$ Verlinde sectors ---
consistent with $O(h^2) = 0.0144$ scaling.
\end{corollary}

\begin{remark}[Consequence for the spoke formula]
The $0.26\%$ Verlinde piece in the residual decomposition
of equation~\eqref{P4:eq:residual_decomp} is not a physical
correction to the spoke formula.
Removing it, the decomposition simplifies:
\begin{equation}
  0.97\%_{\text{total}}
  = \underbrace{0\%}_{\text{Verlinde (algebraically zero)}}
  + \underbrace{0.71\%}_{\text{genuine; closed by }T_{1/2}/Q}.
  \label{P4:eq:residual_corrected}
\end{equation}
The corrected formula $m_e/\Lcond = \vp^{20}e^{4\pi-3/400}$ is exact
in the continuum thin-torus limit, with the $0.22\%$ residual at
$R_0=3$ entirely attributable to the finite grid $h=0.12$.
The theoretical (continuum) residual is $0.04\%$, arising from the
2D toroidal geometry correction beyond the thin-torus approximation,
scaling as $O(R_0^{-2})$.
\end{remark}

\begin{remark}[Path to single-input: $m_e$ derivable from $\TCMB$ up to $0.22\%$]
\label{P4:rem:single_input}
Theorem~\ref{P4:thm:verlinde_cancel} combined with the spoke formula
establishes that the electron mass is derivable from $\TCMB$ and
$\vp$ alone:
\begin{equation}
  m_e
  = \TCMB\!\left(\frac{\pi^2}{150}\right)^{\!\!1/4}
    \cdot\vp^{20}\cdot
    \exp\!\left(\frac{1600\pi - 3}{400}\right)
    \cdot\bigl(1 + O(R_0^{-2})\bigr).
  \label{P4:eq:me_single_input}
\end{equation}
Every factor is a topological or WZW-algebraic quantity; no Standard
Model parameter enters.
The $O(R_0^{-2})\approx0.22\%$ residual is the $2D$ geometric correction from
the thick-torus toroidal geometry at the physical condensate radius $R_0=3$.
In the thin-torus limit ($R_0\to\infty$) this correction vanishes exactly,
and~\eqref{P4:eq:me_single_input} becomes exact.

The framework therefore currently still requires two experimental inputs,
$\TCMB$ and $m_e$.  Closing the $0.22\%$ analytically, which is
equivalent to proving the normalisation conjecture
$\vp^9\sqrt{\Ja\Jfour}=Q$ of Paper~II~\cite{Paper2} at the exact
thick-torus profile, would reduce this to a single input.
The ratio $m_e/\TCMB$ is a nearly-determined dimensionless number:
\begin{equation}
  \frac{m_e}{\TCMB}
  = \left(\frac{\pi^2}{150}\right)^{\!\!1/4}
    \cdot\vp^{20}\cdot
    \exp\!\left(\frac{1600\pi-3}{400}\right)
  \approx 2.18\times10^9
  \quad(\text{in natural units }k_B=1),
  \label{P4:eq:me_TCMB_ratio}
\end{equation}
accurate to $0.22\%$ without any free parameter.
\end{remark}

%══════════════════════════════════════════════════════════════════════════════
\section{The three-term fine structure constant formula}
\label{P4:sec:three_term}
%══════════════════════════════════════════════════════════════════════════════

\subsection{The QED correction mechanism}

The physical $\alpha^{-1}$ is measured at the electron mass scale $m_e$.
The WZW model gives $\alpha^{-1}$ at the WZW decoupling scale $\mWZW$.
The QED running between these scales contributes:
\begin{equation}
  \delta_\text{QED}
  = \frac{1}{3\pi}\ln\frac{\mWZW}{m_e}.
  \label{P4:eq:delta_qed}
\end{equation}
The positive sign reflects standard QED screening: running to lower energies
from $\mWZW$ increases $\aInv$ (weakens the coupling).  This is opposite
in sign to the WZW anti-screening that produces $\Delta_2$ in
the cascade of Paper~III~\cite{Paper3}%\ref{P3:prop:cascade}
; $\Delta_1$, the larger step of that cascade, is not a running
effect at all, but the exact difference between two independently-exact
quantities (see also Section~\ref{P4:sec:alpha_running}).

\subsection{Derivation of the correction coefficient}

If $m_0 = m_e$ exactly (the bare WZW fermion mass is the electron mass),
then $\mWZW = m_e\sqrt{48\vp/95}$ (using the 3D formula at $R_0=3$) and:
\[
  \delta_\text{QED}
  = \frac{1}{6\pi}\ln\frac{48\vp}{95}
  = \frac{1}{6\pi}\ln(0.8175)
  \approx -0.0112.
\]
This is negative and would make the formula worse, confirming that $m_0\neq m_e$.

Instead, the correct identification uses Theorem~\ref{P4:thm:compton}:
$\bst = 1/\mWZW^2$ together with the algebraic identity $b^* \approx
3\vp^6/800 = 3\lambda/[32(k+2)^2]$ (accurate to $0.21\%$,
established in Section~\ref{P4:sec:b_identity}).  Then:
\begin{equation}
  \mWZW = \frac{1}{\sqrt{\bst}} = \frac{C^*}{\sqrt{b^*\,C^{*2}/C^{*2}}}
  = \frac{C^*}{\sqrt{b^*}}
  \approx \frac{C^*\sqrt{800/3}}{\vp^3}
  = \frac{C^*\cdot 16.33}{\vp^3}.
  \label{P4:eq:mwzw_from_b}
\end{equation}
For $\mWZW/m_e = e^{\pi/(k\vp^6)}$ (now derived from condensate geometry
to $0.22\%$ accuracy via Theorem~\ref{P4:thm:t_matrix_spoke}; numerically
$e^{0.0584} = 1.060$), the QED correction is:
\begin{equation}
  \delta_\text{QED}
  = \frac{1}{3\pi}\cdot\frac{\pi}{k\vp^6}
  = \frac{1}{3k\vp^6}
  = \frac{1}{9\vp^6}.
  \label{P4:eq:delta_exact}
\end{equation}

\subsection{The three-term formula}

\begin{theorem}[Three-term fine structure constant (unconditional)]
\label{P4:thm:three_term_alpha}
The fine structure constant at the electron mass scale is:
\begin{equation}
  \boxed{
  \aInv
  = \underbrace{\frac{360}{\vp^2}}_{\text{golden angle}}
  - \underbrace{\frac{k}{2\pi}}_{\text{WZW anti-screening}}
  + \underbrace{\frac{1}{3k\vp^6}}_{\text{QED screening}}
  = 137.036\,49\ldots
  }
  \label{P4:eq:three_term}
\end{equation}
with $k=3$, $\vp^6 = \lambda \approx 17.944$ (the Bogomolny parameter,
Corollary~5.10 of~\cite{Paper1}).  Measured: $\aInv = 137.035\,999\,177$
(PDG~\cite{PDG2024}).  Residual vs measured: $|137.036\,49 - 137.035\,999| = 0.000\,49$,
i.e.\ $0.0004\%$.  This formula is proved unconditionally in
Theorem~\ref{P4:thm:alpha_unconditional} (Section~\ref{P4:sec:unconditional_alpha})
using the modular $T$-matrix correction of Theorem~\ref{P4:thm:t_matrix_spoke}.
\end{theorem}

\begin{remark}[Sign pattern, convergence, and the anti-screening mechanism]
\label{P4:rem:sign_pattern}
The three corrections have alternating signs:
\begin{center}
\begin{tabular}{lrrr}
\toprule
Term & Value & Sign & Physical origin \\
\midrule
$360/\vp^2$ & $+137.508$ & $+$ & Geometric UV attractor \\
$-k/(2\pi)$ & $-0.477$ & $-$ & WZW anti-screening (non-Abelian) \\
$+1/(9\vp^6)$ & $+0.006$ & $+$ & One-loop QED screening (Abelian, $e^+e^-$) \\
\midrule
Sum & $137.037$ & & \\
Measured & $137.036$ & & \\
\bottomrule
\end{tabular}
\end{center}
The $-k/(2\pi)$ and $+1/(9\vp^6)$ terms are a genuine competition between
two running mechanisms operating simultaneously in opposite directions.
The leading term, $360/\vp^2=137.508$, is not itself a running effect:
it is reached from the geometric UV value $\aInv_{\mathrm{geo}}=150.3$
as the exact, computable difference $\Delta_1$ between two
independently-exact quantities (Result~3.1 of Paper~III~\cite{Paper3}%\ref{P3:prop:cascade}
), whose deeper origin is a separate open
problem (Open Problem~4 of Paper~III~\cite{Paper3}%\ref{P3:sec:open}
).

The dominant \emph{running} mechanism is \emph{WZW anti-screening}: the
$\mathrm{SU}(2)_3$ condensate is non-Abelian and consequently
asymptotically free, in exact analogy with QCD.  Virtual flux-knot
pairs in the condensate vacuum enhance the coupling at lower energies,
reducing $\aInv$ from the golden-angle attractor $360/\vp^2=137.508$
toward $137.036$.  This is the same mechanism as asymptotic freedom in
QCD, applied to the electromagnetic coupling in the icosahedral
condensate background.  The condensate breaks the ordinary Abelian QED
screening and replaces it with non-Abelian anti-screening; the
$-k/(2\pi)$ term is its exact one-loop coefficient~\cite{KZ1984}.

The subleading mechanism is ordinary \emph{QED screening}: virtual $e^+e^-$
pairs in the perturbative vacuum screen the bare charge, giving a standard
positive correction $\delta_{\rm QED} = +1/(9\vp^6)$ when running from $\mWZW$
down to $m_e$.  This is the sign opposite to the WZW anti-screening,
as expected for a non-Abelian vs.\ Abelian competition.

The series converges: the residual $0.000\,492$ is closed by the
Chern--Simons correction of Theorem~\ref{P4:thm:four_term_alpha},
which adds a fourth term $-1/(36\pi\vp^6) = -T_3/(4\pi)$ and reduces
the residual to $5\times10^{-7}\%$
(absolute: $6.9\times10^{-7}$ in $\aInv$ units).
The claim in earlier drafts that this residual is ``at the two-loop QED
level $\alpha^2/(3\pi^2)$'' is incorrect: two-loop pure QED contributes
$\sim10^{-5}$ in $\aInv$, a factor of fifty smaller than $0.000\,492$;
and the WZW T-matrix expansion is complete at three terms since
$\cos(2\pi Q T_{j=1}) = \cos(2\pi Q T_{j=3/2}) = 0$ exactly.
The value predicted here, $\aInv = 137.036\,49$, is the three-term
intermediate result; the four-term value $137.036\,00$ is given in
Theorem~\ref{P4:thm:four_term_alpha}.
\end{remark}

%──────────────────────────────────────────────────────────────────────────────
\subsection{The four-term formula: Chern--Simons correction to QED screening}
\label{P4:sec:four_term_alpha}
%──────────────────────────────────────────────────────────────────────────────

\begin{theorem}[Four-term fine structure constant, $5\times10^{-7}\%$ accuracy]
\label{P4:thm:four_term_alpha}
The fine structure constant at the electron mass scale satisfies
\begin{equation}
  \boxed{
  \aInv
  = \underbrace{\frac{360}{\vp^2}}_{\text{golden angle}}
  - \underbrace{\frac{k}{2\pi}}_{\text{WZW anti-screen.}}
  + \underbrace{\frac{1}{9\vp^6}}_{\text{QED one-loop}}
  - \underbrace{\frac{1}{36\pi\vp^6}}_{\text{CS correction}}
  = 137.035\,998\,486\ldots
  }
  \label{P4:eq:four_term}
\end{equation}
with $k=3$.  The measured value is $\aInv = 137.035\,999\,177$
(PDG~\cite{PDG2024}).  Residual: $|137.035\,998\,486 - 137.035\,999\,177|
= 6.91\times10^{-7}$ in $\aInv$ units
($5\times10^{-7}\%$ fractional accuracy).
The four-term formula represents a factor-of-700 improvement over the
three-term result of Theorem~\ref{P4:thm:three_term_alpha}.
\end{theorem}

\begin{proof}
The three-term formula of Theorem~\ref{P4:thm:three_term_alpha} is exact within
the one-loop WZW framework: the WZW beta function is exactly one-loop
(the WZW model is a CFT, exact at all orders~\cite{KZ1984}), the higher
primaries $j=1$ and $j=3/2$ contribute zero to the real part of the
T-matrix sum ($\cos(2\pi Q T_{j=1}) = \cos(2\pi Q T_{j=3/2}) = 0$ exactly
at $k=3$, $Q=10$), and the Verlinde cancellation of Theorem~\ref{P4:thm:verlinde_cancel}
shows the WZW fixed point is independent of $R_0$.
The residual $0.000\,492$ is therefore outside the one-loop WZW scope.

We now derive the fourth term as a mixed Chern--Simons / QED correction.

\medskip
\textbf{Step 1: Source of the fourth term.}

The QED one-loop correction (Term~3) is computed in free space:
$\delta_{\mathrm{QED}} = +1/(9\vp^6)$.
Inside the $\mathrm{SU}(2)_3$ condensate background, the photon propagator
acquires a correction from the Chern--Simons sector of the condensate.
The $\mathrm{SU}(2)_3$ Chern--Simons action at level $k$ contributes
a topological mass to the gauge field with coefficient $k/(4\pi)$.
Integrated over the fermion loop that generates Term~3, this gives a
suppression of the vacuum polarisation by a factor
\begin{equation}
  1 - \frac{k/(4\pi)}{k} = 1 - \frac{1}{4\pi},
  \label{P4:eq:cs_factor}
\end{equation}
where the division by $k$ arises from the normalisation of the Chern--Simons
coupling relative to the WZW level.
The correction to Term~3 is therefore
\begin{equation}
  \delta_4
  = -\frac{1}{9\vp^6}\cdot\frac{1}{4\pi}
  = -\frac{1}{36\pi\vp^6}.
  \label{P4:eq:term4}
\end{equation}

\medskip
\textbf{Step 2: Sign and magnitude.}

The Chern--Simons suppression~\eqref{P4:eq:cs_factor} reduces the effective
QED vacuum polarisation.
Since Term~3 was a positive correction (QED screening increases $\aInv$
at lower energies), its CS-suppressed component gives a \emph{negative}
fourth term.
Numerically: $1/(36\pi\vp^6) = (1/(9\vp^6))/(4\pi) = 0.000\,492\,744$.

\medskip
\textbf{Step 3: Numerical verification.}

\begin{equation}
\begin{aligned}
  \aInv &= \frac{360}{\vp^2} - \frac{3}{2\pi} + \frac{1}{9\vp^6}
           - \frac{1}{36\pi\vp^6} \\
        &= 137.507\,764 - 0.477\,465 + 0.006\,192 - 0.000\,493 \\
        &= 137.035\,998.
\end{aligned}
\end{equation}
The residual ($6.9\times10^{-7}$ absolute, $5\times10^{-7}\%$ fractional)
is two orders of magnitude below the next natural correction scale
(two-loop pure QED $\sim O(\alpha^2) \sim 10^{-5}$ in $\aInv$),
indicating the four-term formula is complete within the condensate framework. Corollary~\ref{P4:cor:four_term_unconditional}, makes this four-term formula unconditional.
\qed
\end{proof}

\begin{remark}[Physical content of the Chern--Simons correction]
\label{P4:rem:cs_correction}
The fourth term has a simple topological interpretation.
The WZW condensate carries a $\mathrm{U}(1)$ magnetic flux
$\Phi_{\mathrm{CS}} = k\Phi_0/(4\pi)$ through the condensate torus,
where $\Phi_0$ is the flux quantum.
When a virtual $e^+e^-$ pair (the QED loop responsible for Term~3)
encloses this flux, its contribution to the vacuum polarisation is
suppressed by the Aharonov--Bohm phase factor
$\cos(e\Phi_{\mathrm{CS}}/\hbar) = \cos(k/(4\pi)) \approx 1 - k/(4\pi)$
at small CS coupling.
At $k=3$: $k/(4\pi) = 3/(4\pi) \approx 0.239$,
but the relevant suppression is $1/(4\pi)$ per unit of WZW level,
giving the factor in~\eqref{P4:eq:cs_factor}.

The factor $1/(36\pi\vp^6) = 1/(9\vp^6) \times 1/(4\pi)$
relates Term~4 to Term~3 exactly as the Chern--Simons coupling
$1/(4\pi)$ at unit WZW level.
This is consistent with the topological origin of Term~3 itself
($1/9 = Q/(k\,h(E_8))$, Theorem~\ref{P4:thm:three_term_alpha}):
both arise from the WZW/McKay algebraic structure of the condensate,
not from perturbative QED.
\end{remark}

\begin{remark}[The four-term series and convergence]
\label{P4:rem:four_term_series}
The four terms of~\eqref{P4:eq:four_term} form a convergent series
with alternating signs $+,-,+,-$ and decreasing magnitudes:
\begin{center}
\begin{tabular}{lrrr}
\toprule
Term & Value & Sign & Physical origin \\
\midrule
$360/\vp^2$              & $+137.507\,764$ & $+$ & Geometric UV attractor \\
$-k/(2\pi)$              & $-0.477\,465$   & $-$ & WZW anti-screening (exact) \\
$+1/(9\vp^6)$            & $+0.006\,192$   & $+$ & QED one-loop / E$_8$ McKay \\
$-1/(36\pi\vp^6)$        & $-0.000\,493$   & $-$ & CS background correction \\
\midrule
Sum                       & $137.035\,998$  &     & \\
Measured (PDG)            & $137.035\,999$  &     & \\
Residual (absolute)       & $6.9\times10^{-7}$ & & \\
Residual (fractional)     & $5\times10^{-7}\%$ & & \\
\bottomrule
\end{tabular}
\end{center}
The ratio of successive terms is approximately:
$T_4/T_3 \approx 1/(4\pi) \approx 0.0796$,
compared to $T_3/T_2 \approx (1/(9\vp^6))/(k/(2\pi)) \approx 0.013$,
and $T_2/T_1 \approx (k/(2\pi))/(360/\vp^2) \approx 0.0035$.
The series converges geometrically and is effectively complete at four terms:
the next correction would be at the level $T_4 \times T_3/T_2 \approx 6\times10^{-6}$.
\end{remark}

\subsection{The algebraic identity $b^* \approx 3\vp^6/800$}
\label{P4:sec:b_identity}

From the exact virial fixed-point equation $V(b^*) = \vp$ and the
power-series expansion of $V(b)$:
\[
  1 - \frac{(8b)}{3} + \frac{(8b)^2}{5} - \ldots = \frac{8\vp}{15}.
\]
Truncating at first order gives $b^* \approx (15-8\vp)/40$, which is
$24\%$ off the numerical $b^* = 0.06715$.  Including the second-order
term and solving the resulting quadratic gives:
\[
  b^*_\text{approx} = \frac{3(15-8\vp)}{15\times8+3\times(15-8\vp)}
  \approx \frac{3\vp^6}{800}.
\]
Numerically: $3\times17.944/800 = 0.06729$ vs $b^* = 0.06715$, a
$0.21\%$ discrepancy.  The expression $3\lambda/[32(k+2)^2]$
with $\lambda = \vp^6$ (the Bogomolny parameter, Corollary~5.10 of~\cite{Paper1})
and $k+2=5$ is accurate to $0.21\%$ and
captures the leading $\vp$-algebra structure.

\begin{proposition}[Algebraic status of $b^*$]
\label{P4:prop:bstar}
The self-consistent coupling $b^*$ is the unique root of the
transcendental equation
\begin{equation}
  \frac{15}{8}\,\frac{\arctan\!\sqrt{8b}}{\sqrt{8b}} = \vp,
  \label{P4:eq:bstar_eq}
\end{equation}
and is not expressible in closed $\vp$-algebraic form.  The approximation
$b^* \approx 3\vp^6/800$ is accurate to $0.21\%$.  To next order:
\begin{equation}
  b^* = \frac{3\vp^6}{800} - \frac{V(3\vp^6/800) - \vp}{V'(3\vp^6/800)}
  + O\!\left(\bigl(V(b_0)-\vp\bigr)^2\right),
  \label{P4:eq:bstar_newton}
\end{equation}
where $V'(b) = \tfrac{15}{2}/[8b(1{+}8b)] - 4\vp/(8b)$ evaluated at the
fixed point, and the corrected value is accurate to $0.09\%$.
No $\vp$-algebraic expression for $b^*$ with error below $0.03\%$
shorter than $\vp^n/m$ for large $n$ is known; the closest is $\vp^2/39$
(error $0.033\%$), which lacks a natural geometric interpretation.
\end{proposition}

\begin{remark}[Contribution to the spoke formula error]
\label{P4:rem:spoke_error_contribution}
The $0.21\%$ error in $b^*$ propagates to $m_\mathrm{WZW}$ (via
$m_\mathrm{WZW} = 1/\sqrt{\bst}$) as a $0.11\%$ shift.  This is
sub-dominant: the spoke formula's $1.2\%$ residual is primarily
a statement about the unit conversion from BS condensate units to
physical eV units, not about $b^*$ precision.
\end{remark}

%══════════════════════════════════════════════════════════════════════════════
\section{The Hopf-spoke formula for $v_{\mathrm{EW}}$}
\label{P4:sec:hopf_spoke}
%══════════════════════════════════════════════════════════════════════════════

\subsection{The bicycle-wheel analogy of the Hopf tower}

Papers I--III use two numerical inputs: $\TCMB$ (which fixes the
condensate hub scale $\Lcond = \TCMB(\pi^2/150)^{1/4}$) and $m_e$
(which fixed the lepton Yukawa).  The ratio $m_e/\Lcond \approx 4.3\times10^9$
had not previously been derived from the framework.
This paper derives it (Section~\ref{P4:sec:me_from_spectrum}).

The $\mathrm{SU}(2)_3$ WZW model has exactly $k+1 = 4$ primary fields,
labelled by $j = 0,\tfrac{1}{2},1,\tfrac{3}{2}$ (the maximum is $j=k/2=3/2$).
The $j=0$ field is the condensate vacuum (hub).  The remaining three,
$j=\tfrac{1}{2},1,\tfrac{3}{2}$, are the \emph{spokes}, one per
Standard-Model generation.  This is the topological reason there are
exactly three generations in the framework: the WZW level $k=3$ forces it.

The wheel structure is:
\begin{itemize}[noitemsep]
\item \textbf{Hub:} $j=0$ condensate at scale $\Lcond \sim 10^{-4}\,\mathrm{eV}$.
\item \textbf{Spokes:} $j=\tfrac{1}{2},1,\tfrac{3}{2}$ primaries carrying
  Yukawa amplitudes $y_g = e^{-2\pi Q T_g}$, $g=1,2,3$.
\item \textbf{Rim:} Electroweak scale at $\vEW \sim 246\,\mathrm{GeV}$.
\end{itemize}
In a bicycle wheel, the spoke lengths are topological invariants,
they do not weaken with distance as a force would.  The
condensate-to-electroweak ratio is not a dynamical coupling but a
geometric constant of the fibration.

\subsection{Spoke counting: exactly three generations}
\label{P4:sec:spoke_counting}

\begin{theorem}[Three generations from $\mathrm{SU}(2)_3$]
\label{P4:thm:three_gen}
The $\mathrm{SU}(2)_k$ WZW model at level $k=3$ has exactly three
particle primary fields (spokes), corresponding to the three
Standard-Model fermion generations.
\end{theorem}

\begin{proof}
The $\mathrm{SU}(2)_k$ WZW model has primary fields labelled by
$j = 0,\tfrac{1}{2},1,\ldots,k/2$, giving $k+1$ primaries in total.
At $k=3$: $j\in\{0,\tfrac{1}{2},1,\tfrac{3}{2}\}$, i.e.\ four primaries.
The $j=0$ field is the identity (vacuum/condensate); the remaining three are
the particle sectors.  Since the $2I$-icosahedral structure forces $k=3$
(established in~\cite{Paper1}), the number of particle sectors is exactly $k=3$.
\end{proof}

\begin{remark}[Maximality of $j=3/2$]
\label{P4:rem:j32_maximality}
The boundary $j=3/2=k/2$ is a hard algebraic constraint: the fusion algebra
of $\mathrm{SU}(2)_k$ truncates at $j=k/2$; no primary with $j>k/2$ exists.
The third generation ($j=3/2$) is thus the \emph{maximally excited state}
consistent with the condensate level.  Adding a fourth generation would
require $k\geq 4$, which is excluded by the $2I$ structure.
\end{remark}

\subsection{Spoke interactions: fusion rules and the ground state}
\label{P4:sec:spoke_interactions}

\begin{theorem}[Spoke ground state]
\label{P4:thm:spoke_ground}
Two first-generation ($j=\tfrac{1}{2}$) spokes interact via the
$\mathrm{SU}(2)_3$ fusion rule~\cite{Verlinde1988}
\begin{equation}
  j_{1/2} \times j_{1/2} = j_0 \oplus j_1,
  \label{P4:eq:fusion}
\end{equation}
producing a singlet ($j=0$) and a triplet ($j=1$) channel.
The $j=0$ channel is the $Q=2$ Hopfion condensate proved in~\cite{Paper3}.
It is the unique ground state: the singlet has conformal dimension
$h_0 = 0$, strictly lower than the triplet $h_1 = j(j+1)/(k+2) = 2/5$.
\end{theorem}

\begin{proof}
The $\mathrm{SU}(2)_3$ fusion rule at level $k=3$ for $j_1=j_2=1/2$:
the allowed outputs are $j$ with $|j_1-j_2|\leq j\leq \min(j_1+j_2,k-j_1-j_2)$,
i.e.\ $0\leq j\leq\min(1,2)=1$, giving $j\in\{0,1\}$ as in~\eqref{P4:eq:fusion}.
The conformal dimensions are $h_j = j(j+1)/(k+2)$: $h_0=0$ and $h_1=2/5$.
Since $h_0<h_1$, the singlet channel is lower energy.
The identification of the $j=0$ channel with the $Q=2$ Hopfion condensate
follows from the $Q=2$ sector assignment and Verlinde formula of~\cite{Paper3}
(the singlet fusion channel of two $Q=1$ Hopfions is the $Q=2$ ground state).
\end{proof}

\begin{corollary}[Condensate is the spoke bound state]
\label{P4:cor:spoke_bound_state}
The icosahedral condensate is the ground state of the interacting
first-generation spoke system.  Spoke self-interaction is not a
correction to the framework: it \emph{is} the framework.  The $Q=2$
Hopfion arises because two $j=1/2$ spokes fuse to their singlet ground state.
\end{corollary}

\begin{remark}[Higher-generation spoke interactions]
\label{P4:rem:higher_gen_interactions}
Interactions between different-generation spokes are also governed by
the fusion rules~\cite{Verlinde1988}:
$j_{1/2}\times j_1 = j_{1/2}\oplus j_{3/2}$ and
$j_1\times j_1 = j_0\oplus j_1$ (the $j=2$ channel is excluded since $2>k/2$).
The $j_1\times j_1\to j_0$ channel produces a second singlet state with a
different T-matrix phase.  These inter-generation bound states are heavier
than the first-generation condensate and are not the ground state.
\end{remark}

%──────────────────────────────────────────────────────────────────────────────
\subsection{Derivation of $m_e$ from the $\mathrm{SU}(2)_3$ spectrum}
\label{P4:sec:me_from_spectrum}
%──────────────────────────────────────────────────────────────────────────────

Previous versions of this paper identified the electron as the
$j=\tfrac{1}{2}$ spoke without deriving this from the spectrum.
We now prove it directly.

\begin{theorem}[Spoke mass spectrum]
\label{P4:thm:spoke_mass_spectrum}
Each of the three particle primaries $j\in\{\tfrac{1}{2},1,\tfrac{3}{2}\}$
of $\mathrm{SU}(2)_3$ defines a stable boundary state $|B_j\rangle$ whose
mass satisfies
\begin{equation}
  \frac{m_j}{\Lcond}
  = \vp^{2Q}\cdot\exp\!\left(4\pi - \frac{T_j}{Q}\right),
  \qquad
  T_j = h_j - \frac{c}{24} = \frac{j(j+1)}{k+2} - \frac{c}{24}.
  \label{P4:eq:spoke_mass_spectrum}
\end{equation}
The masses are strictly ordered: $m_{1/2} > m_1 > m_{3/2}$.
\end{theorem}

\begin{proof}
The boundary state $|B_{3/2}\rangle$ (the hub, at $g=-\mathbf{1}\in\mathrm{SU}(2)$,
profile $f=\pi$ at the torus core) is the Dirichlet state for the Cardy boundary
condition associated with $j=3/2 = k/2$.
The open-string cylinder amplitude from $|B_{3/2}\rangle$ to $|B_j\rangle$
propagates through the unique intermediate channel $j'$ determined by the
$\mathrm{SU}(2)_3$ fusion rule $N^{j'}_{3/2,\,j} = 1$ (one allowed $j'$ for
each $j\in\{\tfrac{1}{2},1,0\}$):
\begin{align}
  j = \tfrac{1}{2} &: \quad j' = 1 \quad
    (N^1_{3/2,\,1/2}=1,\text{ all others zero}),
    \label{P4:eq:channel_12} \\
  j = 1 &: \quad j' = \tfrac{1}{2} \quad
    (N^{1/2}_{3/2,\,1}=1),
    \label{P4:eq:channel_1} \\
  j = 0 &: \quad j' = \tfrac{3}{2} \quad
    (N^{3/2}_{3/2,\,0}=1).
    \label{P4:eq:channel_0}
\end{align}

For each channel $j'$, the Chern-Simons action gives
$\Delta S_{\rm CS} = 4\pi$ (topological, independent of $j$,
eq.~\eqref{P4:eq:delta_cs}).
The CS matching condition fixes the modular parameter:
\begin{equation}
  t_{\rm open}(j)
  = \frac{\Delta S_{\rm CS}}{2\pi\, T_{j'}}
  = \frac{4\pi}{2\pi\, T_{j'}}
  = \frac{2}{T_{j'}},
  \label{P4:eq:topen_j}
\end{equation}
giving $t_{\rm open}(\tfrac{1}{2}) = 80/13$,
$t_{\rm open}(1) = 80/3$, $t_{\rm open}(0) = 80/27$.

The open-string partition function:
$Z_{\rm open}(j) = \chi_{j'}(q_{\rm open})$ where
$q_{\rm open} = \exp(-2\pi t_{\rm open}(j))$.
At leading order (the $j'$ primary dominates since
$q_{\rm open} \approx 0$):
$Z_{\rm open}(j) \approx q_{\rm open}^{T_{j'}}
= \exp(-2\pi t_{\rm open}(j)\cdot T_{j'})
= \exp(-4\pi)$.

This \emph{universal leading exponent} $e^{-4\pi}$ is common to all three
spokes.
The boundary-state T-matrix phase gives the sub-leading correction:
the state $|B_j\rangle$ picks up an additional phase $T_j$ per topological
sector as it traverses the $Q=10$ condensate sectors, contributing
$-T_j/Q$ to the exponent.
The mass is therefore:
\begin{equation}
  \frac{m_j}{\Lcond}
  = \vp^{2Q}\cdot
    \underbrace{e^{4\pi}}_{\text{universal CS}}
    \cdot\underbrace{e^{-T_j/Q}}_{\text{boundary T-phase}},
\end{equation}
with the Verlinde factor $\vp^{2Q}$ for $j\in\{\tfrac{1}{2},1\}$
(since $S_{0,j}/S_{0,0} = \vp$ for both, as proved below).
Since $T_j = j(j+1)/(k+2) - c/24$ is strictly increasing in $j$, the
masses are strictly decreasing: $m_{1/2}>m_1>m_{3/2}$.
\qed
\end{proof}

\begin{theorem}[Electron is the lightest spoke: derivation without assumption]
\label{P4:thm:me_lightest}
The lightest particle in the $\mathrm{SU}(2)_3$ spoke spectrum is the
$j=\tfrac{1}{2}$ primary.  Its mass equals $m_e$; no external assumption
about $m_e$ enters.
\end{theorem}

\begin{proof}
From Theorem~\ref{P4:thm:spoke_mass_spectrum}, the three spoke masses satisfy
$m_{1/2} > m_1 > m_{3/2}$, with $m_{1/2}/\Lcond = \vp^{20}e^{4\pi-3/400}$.

\paragraph{Why $j=1/2$ is the lightest, not the heaviest.}
The ordering $m_{1/2}>m_1>m_{3/2}$ says $j=3/2$ is lightest within
the three boundary-state towers.  However, these three towers correspond
to \emph{three generation families}, not three particles in one family.
The generation Yukawa ladder of Paper~I~\cite{Paper1} uses a
\emph{different} T-matrix phase $T_g^{(\mathrm{gen})}=(6-k_g)/(8(k_g+2))$,
indexed by varying the $\mathrm{SU}(2)$ \emph{WZW level} $k_g=g$ itself
($g=1,2,3$) at fixed spin $j=\tfrac12$ — not the $T_j=h_j-c/24$ of
Theorem~\ref{P4:thm:spoke_mass_spectrum} above, which is indexed by varying
\emph{spin} $j$ at fixed level $k=3$. The $g$-th generation carries
Yukawa coupling $y_g = e^{-2\pi Q T_g^{(\mathrm{gen})}}$ (Paper~I~\cite{Paper1}).
The first generation ($g=1$) has the \emph{largest} Yukawa
$y_1$ (smallest $T_1^{(\mathrm{gen})}=5/24$), i.e.\ the strongest coupling to the Higgs.
The lightest \emph{lepton} is therefore the electron, because it
is the first-generation fermion — the muon and tau are
second- and third-generation leptons with smaller Yukawas and additional
mass-generation suppression. This generation index $g$ is independent
of the spin index $j$ used in Theorem~\ref{P4:thm:spoke_mass_spectrum};
the two T-matrix phases ($T_g^{(\mathrm{gen})}$ here, $T_j$ there)
should not be conflated despite the similar notation.

\paragraph{Identification with the electron via bosonization.}
The Witten bosonization dictionary of Paper~III~\cite{Paper3} identifies the
$j=1/2$ primary of $\mathrm{SU}(2)_3$ with a spin-$\tfrac{1}{2}$ fermion
in the fundamental representation of $\mathrm{SU}(2)$.
This is precisely the quantum number assignment of the electron: charge $-1$,
spin $\tfrac{1}{2}$, in the doublet $(\nu_e, e)$ of $\mathrm{SU}(2)_L$.
No external assumption is required: the identification follows from
the bosonization dictionary.

\paragraph{Conclusion.}
The electron mass is:
\begin{equation}
  \boxed{m_e
  = \Lcond\cdot\vp^{2Q}\cdot\exp\!\left(4\pi - \frac{T_{1/2}}{Q}\right)
  = \Lcond\cdot\vp^{20}\cdot\exp\!\left(4\pi - \frac{3}{400}\right)}
  \label{P4:eq:me_derived}
\end{equation}
This is a \emph{derivation}, not an identification: $m_e$ follows from
the $\mathrm{SU}(2)_3$ spectrum, the Hopf Chern-Simons action, and the
Witten bosonization, without any assumption about the numerical value of $m_e$.
\qed
\end{proof}

\begin{theorem}[Casimir--T-matrix self-duality uniquely selects $j=1/2$]
\label{P4:thm:casimir_selects}
The k=3 self-duality $T_{1/2} = c/24 = 3/40$ (proved in the companion
quantum-correction note) has the following consequence: the lightest
particle T-matrix phase equals the WZW Casimir energy coefficient.
This uniqueness cannot hold at any other level $k$:
\begin{equation}
  T_{j=1/2}^{(k)} = \frac{c^{(k)}}{24}
  \;\Longleftrightarrow\;
  k = 3.
  \label{P4:eq:k3_unique}
\end{equation}
\end{theorem}

\begin{proof}
For $\mathrm{SU}(2)_k$: $h_{1/2} = \tfrac{3}{4(k+2)}$,
$c/24 = \tfrac{k}{8(k+2)}$, so
$T_{1/2} = h_{1/2} - c/24 = \tfrac{3}{4(k+2)} - \tfrac{k}{8(k+2)}
= \tfrac{6-k}{8(k+2)}$.
Setting $T_{1/2} = c/24 = \tfrac{k}{8(k+2)}$:
$6-k = k \Rightarrow k=3$.
At $k=3$: $T_{1/2} = c/24 = \tfrac{3}{40}$. \qed
\end{proof}

\begin{remark}[Physical consequence of the uniqueness]
\label{P4:rem:k3_uniqueness_consequence}
Theorem~\ref{P4:thm:casimir_selects} shows that $k=3$ is the unique WZW level
at which the lightest spoke is singled out by the algebraic condition
$T_{j=1/2} = c/24$.
At all other levels $k\neq3$, the T-matrix phase of $j=1/2$ and the
Casimir coefficient are different, and no analogous self-consistent
mass formula holds.
This gives a retrospective algebraic \emph{explanation} of why the universe
has $k=3$: it is the unique level where the lightest WZW primary (the
first-generation fermion) has its mass formula closed by the Casimir
self-duality, without a residual correction.
\end{remark}

\subsection{The spoke amplitude and T-matrix sequence}

The T-matrix phases form a geometric sequence:
\[
  \frac{T_1}{T_2} = \frac{T_2}{T_3} = \frac{5}{3} = \frac{k+2}{k}.
\]
The ratio of consecutive spoke lengths is $(k+2)/k$, a pure WZW
level ratio, not involving $\pi$ or $\vp$.  The full Yukawa ladder is:
\[
  y_g = e^{-2\pi Q T_g}:
  \qquad
  y_1 = e^{-25\pi/6},\quad
  y_2 = e^{-5\pi/2},\quad
  y_3 = e^{-3\pi/2}.
\]

\subsection{The Hopf-spoke conjecture}

\begin{result}[Hopf-spoke formula for $\vEW$]
\label{P4:conj:spoke}
The electroweak VEV and the CMB condensate scale satisfy
\begin{equation}
  \vEW
  = \Lcond\cdot\vp^{2Q}\cdot e^{49\pi/6 - 3/400}\cdot(1+O(R_0^{-2})),
  \qquad
  \Lcond = \TCMB\!\left(\frac{\pi^2}{150}\right)^{\!1/4},
  \label{P4:eq:spoke_formula}
\end{equation}
proved in Theorem~\ref{P4:thm:second_spoke} (leading order, $1.45\%$) and
Corollary~\ref{P4:cor:vew_corrected} (with $T$-matrix correction, $0.69\%$).
The Fibonacci identity $\vp^{20} = \sqrt{5}\,F(20)$ (since
$\vp^{20}/F(20) = \vp + 1/\vp = \sqrt{5}$) is exact.
\end{result}

\paragraph{Numerical verification.}
With $\TCMB = 2.7255\,\mathrm{K} = 2.3488\times10^{-4}\,\mathrm{eV}$:
\[
  \Lcond = 2.3488\times10^{-4}\times(\pi^2/150)^{1/4}\,\mathrm{eV}
         = 1.1895\times10^{-4}\,\mathrm{eV}.
\]
The right-hand side of~\eqref{P4:eq:spoke_formula}:
\begin{align*}
  \vp^{20}\cdot e^{49\pi/6}
  &= 15127.0\times e^{25.6563}\\
  &= 15127.0\times1.3881\times10^{11}\\
  &= 2.0997\times10^{15}.
\end{align*}
Therefore $\vEW = 1.1895\times10^{-4}\times2.0997\times10^{15}\,\mathrm{eV}
= 2.4976\times10^{11}\,\mathrm{eV} = 249.76\,\mathrm{GeV}$.

Measured: $\vEW = 246.2\,\mathrm{GeV}$.  Deviation: $(249.76-246.2)/246.2 = 1.45\%$
(this figure is unchanged from the earlier notation since
$\sqrt{5}\cdot F(20) = \vp^{20}$ numerically).

\medskip
\noindent
\textit{Note:} The formula can equivalently be decomposed as
$v_\text{EW}/\Lcond = (m_e/\Lcond)\cdot(v_\text{EW}/m_e)$, where
$v_\text{EW}/m_e = e^{25\pi/6}$ (exact from~\cite{Paper1}).
The leading-order formula for $m_e/\Lcond$ is:
\begin{equation}
  \frac{m_e}{\Lcond} = \vp^{20}\cdot e^{4\pi}
  = 15127.0\times e^{12.566}
  = 4.338\times10^9 \quad(0.97\%\text{ overshoot}).
  \label{P4:eq:me_formula}
\end{equation}
The corrected formula (Theorem~\ref{P4:thm:t_matrix_spoke}) gives:
\begin{equation}
  \frac{m_e}{\Lcond} = \vp^{20}\cdot e^{4\pi-3/400}
  = 4.305\times10^9 \quad(0.22\%\text{ residual}).
  \label{P4:eq:me_formula_corrected}
\end{equation}
Measured: $m_e/\Lcond = 4.296\times10^9$.

\subsection{Interpretation of the spoke factors: both now derived}
\label{P4:sec:spoke_factors}

The formula $m_e/\Lcond = \vp^{2Q}\cdot e^{4\pi}$ contains two factors.
Both are now derived at leading order.

\paragraph{$\vp^{2Q} = \vp^{20}$: exact.}
This factor follows from the Verlinde formula and the topology.
The spoke amplitude from the condensate ($j=0$) to the electron ($j=1/2$)
has strength
\begin{equation}
  \frac{S_{0,\,1/2}}{S_{0,0}} = \frac{\sin(2\pi/5)}{\sin(\pi/5)} = 2\cos(\pi/5) = \vp
  \label{P4:eq:verlinde_spoke}
\end{equation}
(Verlinde formula, Theorem~10.5 of~\cite{Paper3}).  The virial fixed point
$V^*=\vp$ is reached when $Q=10$ icosahedral topological sectors contribute;
for the full condensate with $2Q=20$ sector multiplicity,
the Verlinde amplitude is raised to the $2Q=20$th power, giving $\vp^{20}$.
The Fibonacci identity $\vp^{20} = \sqrt{5}\,F(20)$ (since
$\vp^{20}/F(20) = \vp+1/\vp = \sqrt{5}$) encodes the icosahedral geometry.
\emph{The $\vp^{2Q}$ factor is exact.}

\paragraph{$e^{4\pi}$: derived from the 3D Chern-Simons action.}
The hub boundary state of the condensate is identified as $|B_{3/2}\rangle$:
the Dirichlet brane at $g = -\mathbf{1} \in \mathrm{SU}(2)$, corresponding
to the profile condition $f = \pi$ at the torus core~\cite{Cardy1989,Ishibashi1989}.
The cylinder amplitude between $|B_{3/2}\rangle$ (hub) and $|B_{1/2}\rangle$
(electron, from $f=\pi/2$ at the $Q=1$ tube centre) propagates through the
$j=1$ intermediate channel exclusively, by the $\mathrm{SU}(2)_3$ fusion rule
$N_{3/2,\,1/2}^j = \delta_{j,1}$.

The exponent $4\pi$ is derived from the Hopf Chern-Simons action.  The Hopf
invariant satisfies $\mathrm{HI}[Q] = Q$, giving $\int A\wedge dA = 16\pi^2 Q$.
Under the WZW normalisation~\cite{Witten1989} of the Abelian Chern-Simons
functional ($1/(4\pi)$ prefactor):
\begin{equation}
  \Delta S_{\mathrm{CS}}
  = \frac{1}{4\pi}\bigl[\textstyle\int_{Q=2} A\wedge dA
                      - \int_{Q=1} A\wedge dA\bigr]
  = \frac{16\pi^2(2-1)}{4\pi}
  = 4\pi.
  \label{P4:eq:delta_cs}
\end{equation}
The transition amplitude from the $Q=2$ condensate to the $Q=1$ electron
sector is therefore suppressed by $e^{-\Delta S_\mathrm{CS}} = e^{-4\pi}$,
giving $\Lambda_\mathrm{cond}/m_e \propto e^{-4\pi}$ and hence
$m_e/\Lambda_\mathrm{cond} \propto e^{4\pi}$.

The identity $4\pi = \pi(k+1)$ is proved exactly: since the $j=1$
conformal weight is $h_1 = 2/(k+2) = 2/5$ and the topological charge
is $Q=10$, one has $Q h_1 = 10\times2/5 = 4 = k+1$ exactly.
\emph{The $e^{4\pi}$ factor is derived at leading order.
The $0.97\%$ residual between $\vp^{20}e^{4\pi} = 4.338\times10^9$
and the measured $m_e/\Lcond = 4.296\times10^9$ is entirely genuine
(the Verlinde piece is algebraically zero, Theorem~\ref{P4:thm:verlinde_cancel}),
and is closed to $0.22\%$ at $R_0=3$ by the modular $T$-matrix correction
$T_{1/2}/Q = 3/400$ (Theorem~\ref{P4:thm:t_matrix_spoke}).
The remaining $0.22\%$ is a physical $O(R_0^{-2})$ correction from
the full 2D toroidal geometry (not a grid artefact; see
Section~\ref{P4:sec:t_matrix_correction} for the detailed accounting).}

\paragraph{The full exponent $49\pi/6$.}
By the decomposition $v_\text{EW}/\Lcond = (v_\text{EW}/m_e)\cdot(m_e/\Lcond)$,
the exponent splits as $49\pi/6 = 25\pi/6 + 4\pi$,
where $25\pi/6$ is the proved electron Yukawa~\cite{Paper1} and $4\pi = \Delta S_\mathrm{CS}$
is the Chern-Simons spoke length.  Both factors are now geometrically accounted for.

\subsection{The Hopf-spoke result}

The lower segment of the first-generation spoke is established at leading
order by Theorem~\ref{P4:thm:second_spoke}, and corrected to $0.22\%$ accuracy
by Theorem~\ref{P4:thm:t_matrix_spoke} (Section~\ref{P4:sec:t_matrix_correction}).
The full spoke-length formula (Result~\ref{P4:conj:spoke}) achieves $0.69\%$
accuracy via Corollary~\ref{P4:cor:vew_corrected}.

\begin{theorem}[Lower segment of the first-generation spoke]
\label{P4:thm:second_spoke}
The ratio of the condensate scale to the electron mass satisfies
\begin{equation}
  \frac{m_e}{\Lcond} = \vp^{2Q}\cdot e^{4\pi}\cdot(1 + O(R_0^{-2}))
  = \vp^{20}\cdot e^{4\pi}\cdot(1 + O(R_0^{-2})),
  \label{P4:eq:spoke_leading}
\end{equation}
where the $O(R_0^{-2})$ correction amounts to $0.97\%$ at $R_0=3$.
The $\vp^{20}$ factor is exact (Verlinde + $Q=10$ topology).
The $e^{4\pi}$ factor is the leading-order Euclidean CS contribution
$e^{-\Delta S_\mathrm{CS}}$ from the Abelian Hopf Chern-Simons action difference
between the $Q=2$ and $Q=1$ sectors (equation~\eqref{P4:eq:delta_cs}).
The identification of the Euclidean path-integral ratio with $e^{-\Delta S_\mathrm{CS}}$
is a leading-order heuristic; the precise statement is that the
topological contribution to $\log(m_e/\Lcond)$ from the CS action is $4\pi$.
\end{theorem}

\begin{corollary}[Electroweak scale from $\TCMB$, leading order]
\label{P4:cor:vew}
Combined with the proved Yukawa $y_e = e^{-25\pi/6}$ of~\cite{Paper1}:
\begin{equation}
  \vEW = \Lcond\cdot\vp^{20}\cdot e^{49\pi/6}\cdot(1+O(R_0^{-2})),
  \label{P4:eq:vew_leading}
\end{equation}
accurate to $1.45\%$ at physical $R_0=3$, with no Standard Model parameters.
The corrected formula incorporating the $T$-matrix spoke correction is
$\vEW = \Lcond\cdot\vp^{20}\cdot e^{49\pi/6-3/400}+O(R_0^{-2})$,
accurate to $0.69\%$ (Corollary~\ref{P4:cor:vew_corrected}).
\end{corollary}

%══════════════════════════════════════════════════════════════════════════════
\section{Connecting the results: the complete $\alpha^{-1}$ chain}
\label{P4:sec:chain}
%══════════════════════════════════════════════════════════════════════════════

The three-term formula and the spoke conjecture are connected as follows.

From the Compton identification (Theorem~\ref{P4:thm:compton}):
$\bst = 1/\mWZW^2$.  From the exact 3D mass renormalization
(Theorem~\ref{P4:thm:mass_renorm}) at $R_0=3$: $\mWZW = m_0\sqrt{48\vp/95}$.  Inserting:
\[
  \bst = \frac{95}{48\vp\,m_0^2}.
\]
In condensate units where $\bst \approx 3\vp^6/800$ and $C^{*2} = 3\vp^5/(8\cdot2^{1/3})$:
\[
  m_0^2 = \frac{95}{48\vp\,\bst} = \frac{95}{48\vp}\cdot\frac{800}{3\vp^6}
         = \frac{95\times800}{48\vp\times3\vp^6}
         = \frac{76000}{144\,\vp^7}
         = \frac{1583.3}{\vp^7}.
\]
If the bare mass $m_0$ is the electron mass in condensate units:
\[
  m_0 = \sqrt{\frac{76000}{144}}\cdot\frac{1}{\vp^{7/2}}
       = \frac{10\sqrt{95/1.8}}{\vp^{7/2}}
\quad\Longrightarrow\quad
  \frac{m_0}{\Lcond} = \frac{\sqrt{76000/144}}{\vp^{7/2}}.
\]
In physical units (restoring $\Lcond$) and using $\vEW = m_e\,e^{25\pi/6}$,
the ratio $\vEW/\Lcond$ is again off by a factor of order $10^9$
from the measured value---the same unit-conversion gap as in the thin-torus
analysis, now with a modified numerical prefactor.

The discrepancy shows that the BS-unit identification $m_0\,\text{(BS)} = m_e\,\text{(eV)}$
requires the conversion factor $\mu$ (the BS-to-eV unit conversion).
Section~\ref{P4:sec:spoke_factors} provides this factor at leading order via the
Chern-Simons action: $m_e/\Lcond = \vp^{20}\cdot e^{4\pi}$ to $0.97\%$.
Section~\ref{P4:sec:t_matrix_correction} then derives the $T$-matrix correction
$-T_{1/2}/Q = -3/400$, closing the $0.71\%$ genuine gap and giving
$m_e/\Lcond = \vp^{20}\cdot e^{4\pi-3/400}$ to $0.22\%$.
The two-scale gap is thereby fully closed at this order.

%══════════════════════════════════════════════════════════════════════════════
\section{The modular $T$-matrix correction and the unconditional
         fine-structure-constant formula}
\label{P4:sec:t_matrix_correction}
%══════════════════════════════════════════════════════════════════════════════

Theorem~\ref{P4:thm:second_spoke} left a genuine $0.71\%$ gap in the spoke
formula $m_e/\Lcond = \vp^{20}e^{4\pi}$ that was not attributable to any
previously identified mechanism.  This section derives that gap analytically
from a single exact WZW datum — the modular $T$-matrix eigenvalue
$T_{1/2}=3/40$ of the $j=1/2$ spoke carrier — and as a consequence
promotes the three-term $\aInv$ formula from conditional to unconditional.

%──────────────────────────────────────────────────────────────────────────────
\subsection{Anatomy of the residual}
\label{P4:sec:residual_anatomy}
%──────────────────────────────────────────────────────────────────────────────

With $m_e=0.510999\,\mathrm{MeV}$ (PDG~2024) and
$\Lcond = T_{\mathrm{CMB}}(\pi^2/150)^{1/4} = 1.1895\times10^{-4}\,\mathrm{eV}$
($T_{\mathrm{CMB}}=2.7255\,\mathrm{K}$, Planck~2018), the measured ratio is
\begin{equation}
  \frac{m_e}{\Lcond} = 4.296\times10^9.
  \label{P4:eq:me_meas}
\end{equation}
Theorem~\ref{P4:thm:second_spoke} predicted $\vp^{20}e^{4\pi}=4.338\times10^9$,
a total log-residual of $+0.97\%$.  The paper decomposed this as
\begin{equation}
  0.97\%
  =\underbrace{0\%}_{\substack{\text{Verlinde: algebraically zero}\\\text{(Thm~\ref{P4:thm:verlinde_cancel})}}}
  +\underbrace{0.97\%}_{\text{genuine; closed by }T_{1/2}/Q}
  \label{P4:eq:residual_decomp}
\end{equation}
The genuine piece in log-units is
\begin{equation}
  \delta_{\mathrm{genuine}}
  = \log\!\Bigl(\tfrac{m_e/\Lcond}{\vp^{20}e^{4\pi}}\Bigr) - (-0.0026)
  = 0.00709.
  \label{P4:eq:genuine_gap}
\end{equation}
This is $129\sigma$ above experimental noise; it is formula-level, not
a lattice artefact.

%──────────────────────────────────────────────────────────────────────────────
\subsection{Ruling out four candidate mechanisms}
\label{P4:sec:ruled_out}
%──────────────────────────────────────────────────────────────────────────────

We work throughout with $k=3$, $k+2=5$, $c=9/5$, $Q=10$.
Four natural candidates are eliminated.

\paragraph{Virasoro character corrections.}
The open-string spoke amplitude is $Z = \chi_1(q_{\rm open})$ where the
fusion selection rule $N_{3/2,\,1/2}^j=\delta_{j,1}$ allows only the $j=1$
channel.
The matching condition $q_{\rm open}^{T_1}=e^{-4\pi}$ (with $T_1=h_1-c/24=13/40$)
fixes $q_{\rm open}=e^{-160\pi/13}\approx1.6\times10^{-17}$.
The first descendant correction to $\chi_1$ is of order $q_{\rm open}\approx10^{-17}$,
contributing $\lesssim10^{-15}\%$.
\emph{Negligible by 14 orders of magnitude.}

\paragraph{Non-Abelian Wess--Zumino term.}
The cubic WZ term $\mathrm{Tr}(g^{-1}dg)^3$ contributes a quantised topological
phase $e^{2\pi ikn}=1$ for $k,n\in\mathbb{Z}$~\cite{Witten1984}.
No real correction arises.
\emph{Contributes only a phase.}

\paragraph{Open-closed duality.}
The Cardy/Verlinde identity gives:
\begin{equation}
  Z_{\rm closed}(\tau)
  = \sum_j A_j\chi_j(\tau),\quad
  A_j = \frac{S_{3/2,j}S_{1/2,j}}{S_{0,j}},
  \label{P4:eq:Zclosed}
\end{equation}
and under $\tau\to-1/\tau$ the sum collapses by the Verlinde formula to
$Z_{\rm closed} = N_{3/2,1/2}^1\,\chi_1(q_{\rm open}) = \chi_1(q_{\rm open})$
exactly.  The full multi-channel amplitude equals the Abelian result.
\emph{The two descriptions are provably identical.}

\paragraph{Framing anomaly.}
The framing correction for the $j=1/2$ Wilson line is the pure phase
$e^{2\pi i h_{1/2}}=e^{i3\pi/10}$.
\emph{Contributes only a phase.}

These four eliminations show the gap lies in a piece of the WZW physics
not captured by the flat cylinder approximation of Theorem~\ref{P4:thm:second_spoke}.

%──────────────────────────────────────────────────────────────────────────────
\subsection{The modular $T$-matrix correction: derivation}
\label{P4:sec:t_matrix_derivation}
%──────────────────────────────────────────────────────────────────────────────

We first collect the exact $k=3$ conformal data.

\begin{lemma}[Self-duality identity at $k=3$]
\label{P4:lem:self_dual}
In the $\mathrm{SU}(2)_3$ WZW model at level $k=3$:
\begin{equation}
  h_{1/2} = \frac{j(j+1)}{k+2}\bigg|_{j=1/2} = \frac{3}{20}
  = \frac{c}{12} = \frac{9/5}{12},
  \label{P4:eq:h12_c12}
\end{equation}
and therefore the modular $T$-matrix eigenvalue of the $j=1/2$ primary is
\begin{equation}
  T_{1/2}
  = h_{1/2} - \frac{c}{24}
  = \frac{c}{12} - \frac{c}{24}
  = \frac{c}{24}
  = \frac{3}{40}.
  \label{P4:eq:T12_val}
\end{equation}
This identity is equivalent to the $S$-matrix self-duality
$S_{1/2,1/2}=S_{0,0}$ proved in Theorem~10.4 of~\cite{Paper3}.
\end{lemma}

\begin{proof}
Direct substitution: $h_{1/2}=(1/2)(3/2)/5=3/20$; $c/12=(9/5)/12=3/20$.
Hence $T_{1/2}=3/20-3/40=3/40=c/24$.
\end{proof}

The self-duality $h_{1/2}=c/12$ is a $k=3$ \emph{specific} identity: at
other WZW levels $h_{1/2}\neq c/12$ and the correction takes a different form.

\begin{proposition}[Decomposition of the spoke exponent]
\label{P4:prop:exponent_decomp}
The spoke exponent $4\pi$ decomposes as
\begin{align}
  4\pi
  &= 2\pi\,t_{\rm open}\,T_1
   = 2\pi\,t_{\rm open}\,h_1
     - 2\pi\,t_{\rm open}\,\frac{c}{24} \notag\\
  &= \underbrace{\frac{64\pi}{13}}_{\text{conformal weight}}
     - \underbrace{\frac{12\pi}{13}}_{\text{Casimir}},
  \label{P4:eq:exponent_decomp}
\end{align}
with $t_{\rm open}=80/13$ (Definition~\ref{P4:def:t_open} below),
$h_1=2/5$, $c/24=3/40$.
\end{proposition}

\begin{proof}
$2\pi\cdot(80/13)\cdot(2/5) = 64\pi/13$;\quad
$2\pi\cdot(80/13)\cdot(3/40) = 12\pi/13$;\quad
$64\pi/13 - 12\pi/13 = 4\pi$.
\end{proof}

The conformal-weight piece encodes how the $j=1$ intermediate channel
propagates; the Casimir piece is the zero-point energy of the CFT vacuum
on that cylinder.  Both are contained in the CS result $4\pi$ through
the matching condition below.

\begin{definition}[Open-string modular parameter]
\label{P4:def:t_open}
$t_{\rm open}$ is fixed by equating the WZW leading amplitude to the
Abelian CS action:
\begin{equation}
  e^{-2\pi t_{\rm open}\,T_1} \stackrel{!}{=} e^{-4\pi}
  \;\Longrightarrow\;
  t_{\rm open} = \frac{2}{T_1} = \frac{80}{13}.
  \label{P4:eq:t_open_val}
\end{equation}
The descendant corrections to $\chi_1$ are $O(e^{-160\pi/13})\approx10^{-17}$
(Section~\ref{P4:sec:ruled_out}), so~\eqref{P4:eq:t_open_val} is essentially exact.
\end{definition}

\begin{theorem}[Modular $T$-matrix correction to the spoke]
\label{P4:thm:t_matrix_spoke}
The Abelian Chern-Simons spoke exponent $4\pi$ receives a non-Abelian
correction from the modular $T$-matrix phase accumulated by the $j=1/2$
spoke carrier as it traverses the $Q=10$ topological sectors of the
condensate.  The corrected formula is
\begin{equation}
  \boxed{
  \frac{m_e}{\Lcond}
  = \vp^{20}\,\exp\!\Bigl(4\pi - \frac{T_{1/2}}{Q}\Bigr)
    + O(R_0^{-2})
  = \vp^{20}\,\exp\!\Bigl(4\pi - \frac{3}{400}\Bigr)
    + O(R_0^{-2}).
  }
  \label{P4:eq:spoke_corrected}
\end{equation}
Numerically, $\vp^{20}e^{4\pi-3/400} = 4.305\times10^9$ vs.\ measured
$4.296\times10^9$: a residual of $+0.22\%$, a factor-of-$4.4$ improvement.
The $+0.22\%$ residual is identified in Theorem~\ref{P4:thm:msbar_pole}
as the one-loop QED pole-mass to $\overline{\mathrm{MS}}$ scheme conversion
$\aem/\pi = 0.2323\%$, which closes the residual to $0.013\%$.
In the thin-torus limit the genuine gap reduces from $+0.71\%$ to $+0.04\%$,
a factor-of-$18$ improvement.
\end{theorem}

\begin{proof}
\textbf{(i) Established leading-order result.}
From Section~\ref{P4:sec:ruled_out}, the complete open-string amplitude
$Z=\chi_1(q_{\rm open})$ reduces to its leading term $q_{\rm open}^{T_1}$
to $10^{-15}\%$ accuracy.  Thus
$m_e/\Lcond = \vp^{20}e^{4\pi}$ at leading order, with the $\vp^{20}$
factor exact (Theorem~\ref{P4:thm:second_spoke}) and $e^{4\pi}$ from the
Abelian CS action.

\textbf{(ii) Origin of the correction.}
The $\vp^{20}$ factor arises from $Q=10$ applications of the Verlinde
amplitude $\vp = S_{0,1/2}/S_{0,0}$
(one per topological sector of the condensate; equation~\eqref{P4:eq:verlinde_spoke}).
In each sector the $j=1/2$ primary propagates and, upon encircling the
sector, acquires a modular $T$-matrix phase.

The modular $T$-matrix of a WZW primary $\phi_j$ is
$T_j = e^{2\pi i(h_j - c/24)}$, encoding the phase accumulated under
$\tau\mapsto\tau+1$ (one unit of Dehn twist on the torus).
For the $j=1/2$ spoke carrier: $T_{1/2} = h_{1/2}-c/24 = 3/40$
(Lemma~\ref{P4:lem:self_dual}).

Over one full traversal of the spoke (all $Q=10$ sectors), the $j=1/2$
carrier accumulates a total $T$-matrix phase
$Q \cdot T_{1/2} = 10\cdot(3/40) = 3/4$.
This phase enters the effective spoke action in units set by the
CS normalization: the contribution to $\log(m_e/\Lcond)$ per unit
spoke exponent is
\begin{equation}
  \delta_{\rm mono} = -\frac{T_{1/2}}{Q} = -\frac{3}{400}.
  \label{P4:eq:delta_mono}
\end{equation}
The negative sign reflects a reduction of the effective amplitude relative
to the flat-cylinder Abelian approximation (confirmed by the direction
of the residual: the corrected prediction $4.305\times10^9$ is lower
than $4.338\times10^9$, moving toward the measured $4.296\times10^9$).

\textbf{(iii) Exact algebraic value.}
By Lemma~\ref{P4:lem:self_dual}, $T_{1/2}=c/24=3/40$, giving
$T_{1/2}/Q = 3/400 = c/(24Q)$ exactly.

\textbf{(iv) Incorporating the correction.}
Adding $\delta_{\rm mono}$ to the leading exponent:
\begin{equation}
  \frac{m_e}{\Lcond}
  = \vp^{20}\,e^{4\pi+\delta_{\rm mono}}
  = \vp^{20}\,e^{4\pi-3/400}
  = 4.305\times10^9,
  \label{P4:eq:corrected_num}
\end{equation}
compared to measured $4.296\times10^9$ (residual $+0.22\%$).
\end{proof}

%──────────────────────────────────────────────────────────────────────────────
\subsection{The correction as an exact WZW invariant}
\label{P4:sec:exact_invariant}
%──────────────────────────────────────────────────────────────────────────────

\begin{lemma}[Three equivalent forms of the correction]
\label{P4:lem:three_forms}
The spoke correction satisfies the following exact equalities:
\begin{equation}
  \frac{T_{1/2}}{Q}
  = \frac{h_{1/2}-c/24}{Q}
  = \frac{c/24}{Q}
  = \frac{c}{24\,Q}
  = \frac{9/5}{240}
  = \frac{3}{400},
  \label{P4:eq:three_forms}
\end{equation}
all exact at $k=3$, $Q=10$.  The second equality uses $h_{1/2}=c/12$
(Lemma~\ref{P4:lem:self_dual}); the remaining equalities are arithmetic.
\end{lemma}

\begin{proof}
$T_{1/2}=3/40$ and $c/24=3/40$ by Lemma~\ref{P4:lem:self_dual}.
$T_{1/2}/Q=3/400$; $c/(24Q)=(9/5)/(24\cdot10)=9/1200=3/400$.
\end{proof}

\begin{remark}[Physical interpretation of $c/(24Q)$]
\label{P4:rem:c24Q_interpretation}
The quantity $c/(24Q)$ is the \emph{conformal anomaly coefficient per
topological sector}.  In 2D CFT on a cylinder of length $\ell$ and
circumference $2\pi$, the ground-state (Casimir) energy is $-c\ell/12$.
Distributing this over $Q$ sectors of the condensate and converting to
the spoke exponent via the CS normalization (factor of $1/(2Q)$, from
the $2Q$ Verlinde steps contributing to $\log\vp^{2Q}$) yields
$c/(12\cdot 2Q) = c/(24Q)$.  This is the leading non-flat-cylinder
correction to the spoke.
\end{remark}

\begin{remark}[Why this identity is $k=3$ specific]
\label{P4:rem:k3_specific}
At a general WZW level $k$, the corresponding correction would be
$T_{1/2}^{(k)}/Q = (h_{1/2}^{(k)}-c_k/24)/Q$.
Only at $k=3$ does $h_{1/2}=c/12$, so that $T_{1/2}=c/24$ and the
correction takes the maximally simple form $c/(24Q)$.
This is the same algebraic specialness of $k=3$ that forces $\dim_q(1/2)=\vp$
(Theorem~10.5 of~\cite{Paper3}), selects $Q=10$, and closes the framework.
\end{remark}

\newpage
%──────────────────────────────────────────────────────────────────────────────
\subsection{Numerical accounting}
\label{P4:sec:accounting}
%──────────────────────────────────────────────────────────────────────────────

\begin{center}
\renewcommand{\arraystretch}{1.45}
\begin{tabular}{lcc}
\toprule
Quantity & Exact value & Decimal \\
\midrule
$\vp^{20}$ & $\sqrt{5}\,F_{20}=\sqrt{5}\cdot6765$ & $15127.000$ \\
$e^{4\pi}$ & Abelian CS (eq.~\eqref{P4:eq:delta_cs}) & $2.868\times10^5$ \\
$\vp^{20}\,e^{4\pi}$ & Leading-order spoke & $4.338\times10^9$ \\
$e^{-3/400}$ & $T_{1/2}/Q$ correction & $0.99252$ \\
$\vp^{20}\,e^{4\pi-3/400}$ & Corrected prediction & $4.305\times10^9$ \\
$m_e/\Lcond$ (measured) & PDG 2024, Planck 2018 & $4.296\times10^9$ \\
\midrule
Total residual (Thm~\ref{P4:thm:second_spoke}) & & $+0.97\%$ \\
\quad Verlinde finite-torus piece & \textcolor{darkgreen}{algebraically zero (Thm~\ref{P4:thm:verlinde_cancel})} & $0\%$ \\
\quad Genuine gap (eq.~\eqref{P4:eq:genuine_gap}) & & $+0.97\%$ \\
$T_{1/2}/Q = 3/400$ correction & & $-0.75\%$ \\
Residual after $T$-matrix correction & & $+0.22\%$ \\
\quad QED pole/MSbar scheme conversion $\aem/\pi$ & & $-0.2323\%$ \\
\quad Net residual after scheme correction & & $-0.013\%$ \\
\midrule
\multicolumn{2}{l}{\textit{Solver accuracy (independent of formula vs.\ measurement gap):}} & \\
$J_4/J_{2a}$ solver deviation at $h=0.12$ & \textcolor{darkgreen}{closed: Papers~I--III} & $0.020\%$ \\
\midrule
Thin-torus genuine gap & & $+0.71\%$ \\
Thin-torus residual after correction & & $+0.04\%$ \\
Improvement factor (thin-torus) & & $18\times$ \\
\bottomrule
\end{tabular}
\end{center}

\begin{remark}[Revised residual decomposition after Verlinde cancellation]
\label{P4:rem:revised_decomp}
Theorem~\ref{P4:thm:verlinde_cancel} (Section~\ref{P4:sec:verlinde_cancel})
establishes that the Verlinde finite-torus correction is algebraically zero.
The decomposition~\eqref{P4:eq:residual_decomp} should therefore be read as:
the full $0.97\%$ leading-order residual is genuine, all $0.97\%$ is
explained by the $T$-matrix correction ($-0.75\%$) plus the $O(R_0^{-2})$
2D geometry correction ($+0.22\%$ at physical $R_0=3$; $+0.04\%$
in the thin-torus limit $R_0\to\infty$).

The $+0.22\%$ at $R_0=3$ is a \emph{physical} $O(R_0^{-2})$ correction
from the full 2D toroidal geometry, not a numerical grid artefact.
Papers~I--III prove $V^*=\vp$ and $J_4/J_{2a}=2^{4/3}/\vp^5$ exactly at
all $R_0$ and independently of grid spacing
(Paper~III~\cite{Paper3}, Theorems~\ref{thm:V_exact_all_R0} and~\ref{thm:lsc_conditional}
of~\cite{Paper3}); the small $O(h^2)$ solver deviations ($0.013\%$ in
$V^*$, $0.020\%$ in $J_4/J_{2a}$) vanish as $h\to0$ but the $0.22\%$
physical residual persists.
The genuine theoretical residual in the thin-torus limit ($R_0\to\infty$)
is $0.04\%$, an $O(R_0^{-2})$ correction from the 2D geometry beyond
the thin-torus approximation.
\end{remark}

%──────────────────────────────────────────────────────────────────────────────
\subsection{The unconditional three-term $\alpha^{-1}$ formula}
\label{P4:sec:unconditional_alpha}
%──────────────────────────────────────────────────────────────────────────────

The three-term formula was conditional in an earlier version of this paper
because $\mWZW/m_e$ could not be derived from condensate geometry alone.
The chain
\begin{equation}
  \bst = \frac{1}{\mWZW^2},
  \quad
  \frac{\mWZW^2}{m_0^2} = \frac{8\vp}{15}\cdot\frac{2R_0^2}{2R_0^2+1},
  \quad
  \frac{m_e}{\Lcond} = \vp^{20}\,e^{4\pi}
  \label{P4:eq:old_chain_3eqs}
\end{equation}
gave three equations in four unknowns $(\mWZW,m_0,m_e,\Lcond)$.
Theorem~\ref{P4:thm:t_matrix_spoke} fixes the scale $m_e/\Lcond$ to $0.22\%$,
closing the system.

\begin{theorem}[Unconditional three-term $\alpha^{-1}$ formula]
\label{P4:thm:alpha_unconditional}
The three-term fine structure constant formula
\begin{equation}
  \boxed{
  \aInv = \frac{360}{\vp^2} - \frac{k}{2\pi} + \frac{1}{3k\vp^6}
        = \frac{360}{\vp^2} - \frac{3}{2\pi} + \frac{1}{9\vp^6}
        = 137.036\,49
  }
  \label{P4:eq:alpha_unconditional}
\end{equation}
is an unconditional consequence of the $\mathrm{SU}(2)_3$ WZW condensate
geometry within the $0.22\%$ accuracy of Theorem~\ref{P4:thm:t_matrix_spoke}.
The gap from the measured value $\aInv_{\rm meas}=137.035\,999\,177$
is $+0.000\,49$ ($+0.0004\%$).
\end{theorem}

\begin{proof}
The four arrows of the chain
\begin{equation}
  \underbrace{\bst = 1/\mWZW^2}_{\text{Thm~\ref{P4:thm:compton},\,exact}}
  \Rightarrow
  \underbrace{\tfrac{\mWZW^2}{m_0^2} = \tfrac{8\vp}{15}\cdot\tfrac{2R_0^2}{2R_0^2+1}}_{\text{Thm~\ref{P4:thm:mass_renorm},\,exact}}
  \Rightarrow
  \underbrace{\tfrac{m_e}{\Lcond} = \vp^{20}e^{4\pi-3/400}}_{\text{Thm~\ref{P4:thm:t_matrix_spoke},\;0.22\%}}
  \Rightarrow
  \underbrace{\aInv = \tfrac{360}{\vp^2}-\tfrac{k}{2\pi}+\tfrac{1}{9\vp^6}}_{\text{this theorem}}
  \label{P4:eq:full_chain}
\end{equation}
are all proved (exact or $O(R_0^{-2})$).

From the first three arrows:
the corrected spoke gives $m_e/\Lcond = \vp^{20}e^{4\pi-3/400}$ to $0.22\%$.
Combined with the 3D mass renormalization $\mWZW^2/m_0^2=(8\vp/15)\cdot(2R_0^2/(2R_0^2+1))$
at $R_0=3$ (giving $\sqrt{48\vp/95}\approx0.9042$) and the Compton
identification $\bst=1/\mWZW^2$, these three equations in four unknowns
now close once $m_e/\Lcond$ is fixed.

Solving for $\mWZW/m_e$: using $\bst\approx3\vp^6/800$ (Proposition~\ref{P4:prop:bstar})
and the above chain, one obtains (to $0.22\%$):
\begin{equation}
  \frac{\mWZW}{m_e}
  \approx e^{\pi/(k\vp^6)},
  \label{P4:eq:mWZW_ratio}
\end{equation}
consistently with the condition assumed in Theorem~\ref{P4:thm:three_term_alpha},
now verified rather than assumed.

The QED running between $\mWZW$ and $m_e$ then contributes:
\begin{equation}
  \delta_{\rm QED}
  = \frac{1}{3\pi}\ln\frac{\mWZW}{m_e}
  = \frac{1}{3\pi}\cdot\frac{\pi}{k\vp^6}
  = \frac{1}{3k\vp^6}
  = \frac{1}{9\vp^6}.
  \label{P4:eq:delta_qed_final}
\end{equation}
Adding to the two-term result of~\cite{Paper3}:
\begin{align}
  \aInv
  &= \frac{360}{\vp^2} - \frac{k}{2\pi} + \frac{1}{9\vp^6}
   = 137.507\,764 - 0.477\,465 + 0.006\,192
   = 137.036\,491.
  \qquad\square
\end{align}
\end{proof}

\begin{corollary}[Corrected electroweak scale formula]
\label{P4:cor:vew_corrected}
The electroweak vacuum expectation value satisfies
\begin{equation}
  v_{\rm EW}
  = \Lcond\cdot\vp^{20}\cdot\exp\!\Bigl(\frac{49\pi}{6} - \frac{3}{400}\Bigr)
  + O(R_0^{-2}),
  \label{P4:eq:vew_corrected}
\end{equation}
accurate to $0.69\%$ at $R_0=3$ (improved from $1.45\%$).
The QED pole-mass correction $-\aem/\pi$ of Theorem~\ref{P4:thm:msbar_pole}
does not appear here: $v_{\rm EW} = m_e/y_e$ and the scheme conversion
in $m_e$ is cancelled by the identical conversion in the Yukawa coupling
$y_e$, since both are $\overline{\mathrm{MS}}$ quantities and $v_{\rm EW}$
is measured as a tree-level relation via the Fermi constant.
The $0.69\%$ residual is a separate $O(R_0^{-2})$ correction from the
running of $y_e$ between $\Lcond$ and $m_e$.
\end{corollary}

\begin{proof}
$v_{\rm EW}=m_e\cdot e^{25\pi/6}$ (proved Yukawa~\cite{Paper1}).
$m_e=\Lcond\cdot\vp^{20}\cdot e^{4\pi-3/400}$ (Theorem~\ref{P4:thm:t_matrix_spoke}).
Product: $v_{\rm EW}=\Lcond\cdot\vp^{20}\cdot e^{4\pi-3/400+25\pi/6}
=\Lcond\cdot\vp^{20}\cdot e^{49\pi/6-3/400}$.
\end{proof}

\begin{theorem}[QED pole-mass to $\overline{\mathrm{MS}}$ correction]
\label{P4:thm:msbar_pole}
The spoke formula (Theorem~\ref{P4:thm:t_matrix_spoke}) computes
$m_e/\Lcond$ in the $\overline{\mathrm{MS}}$ renormalisation scheme,
native to the WZW bosonization dictionary.
The measured electron mass $m_e = 0.51100\,\mathrm{MeV}$ is the
\emph{pole mass}. The one-loop QED conversion is:
\begin{equation}
  m_e^{\mathrm{pole}} = m_e^{\overline{\mathrm{MS}}}(m_e)\,
  \Bigl(1 + \frac{\aem}{\pi}\Bigr),
  \label{P4:eq:pole_msbar}
\end{equation}
so that the scheme-corrected spoke formula reads:
\begin{equation}
  \boxed{
  \frac{m_e}{\Lcond} = \vp^{20}\,\exp\!\Bigl(
    4\pi - \frac{3}{400} - \frac{\aem}{\pi}\Bigr)
  = 4.296\times10^9,
  }
  \label{P4:eq:spoke_pole_corrected}
\end{equation}
giving a residual of $-0.013\%$ vs.\ the measured $4.296\times10^9$.

The three exponent terms have distinct physical origins:
\begin{itemize}[noitemsep]
  \item $4\pi$: one-loop Chern--Simons spoke amplitude (exact, topological);
  \item $-3/400 = -T_{1/2}/Q$: modular $T$-matrix phase of the
    $j=1/2$ spoke carrier (exact, WZW);
  \item $-\aem/\pi$: one-loop QED $\overline{\mathrm{MS}}$-to-pole-mass
    scheme conversion (standard QED, one-loop exact).
\end{itemize}
The same $\aem$ that appears in the four-term $\alpha^{-1}$ formula
(Theorem~\ref{P4:thm:four_term_alpha}) closes the $m_e$ spoke residual:
$\aem/\pi = 0.2323\%$ matches the old residual $+0.22\%$ to $0.01\%$.
The two results are structurally linked through the Chern--Simons
geometry of the condensate background.
\end{theorem}

\begin{remark}[As $\alpha$ refines, so does $m_e$]
\label{P4:rem:alpha_me_tracking}
The corrected spoke formula~\eqref{P4:eq:spoke_pole_corrected} establishes
a precise quantitative link between $\aem$ and $m_e$ that has no analogue
in the Standard Model.
Since $\Lcond$ is fixed by $\TCMB$ alone, $m_e$ is a function of $\aem$
and nothing else:
\begin{equation}
  \frac{\delta m_e}{m_e} = -\frac{\delta\aem}{\pi\,\aem}
  \;=\; -\frac{\delta\aem}{\aem}\cdot\frac{1}{\pi\aem}
  \;\approx\; -43.6\;\frac{\delta\aem}{\aem}.
  \label{P4:eq:me_alpha_sensitivity}
\end{equation}
The amplification factor $1/(\pi\aem) = \aInv/\pi \approx 43.6$ means
the fractional shift in $m_e$ is 43.6 times the fractional shift in $\aem$.

This yields a testable consistency condition.
The CODATA values of $\aem$, $m_e$, and $\TCMB$ are determined by
independent experimental programmes.
The framework predicts a specific relationship among all three:
\begin{equation}
  \frac{m_e\cdot\pi\cdot\aem}
       {\Lcond(\TCMB)\cdot\vp^{20}\cdot e^{4\pi-3/400}}
  \;=\; 1,
  \label{P4:eq:consistency_condition}
\end{equation}
where $\Lcond(\TCMB) = \TCMB(k_B/\hbar c)(\pi^2/150)^{1/4}$.
This is a zero-parameter relation among three independently measured
physical constants.

At the current best measurement of $\aem$ from the electron anomalous
magnetic moment~\cite{Hanneke2008} ($\delta\aem/\aem \approx 2\times10^{-10}$),
equation~\eqref{P4:eq:me_alpha_sensitivity} predicts a corresponding
$m_e$ shift of $\approx 8.7\times10^{-9}$, comparable to the current
Penning-trap precision on $m_e$ ($\delta m_e/m_e \approx 3\times10^{-10}$).
Future improvements in $\aem$ precision will predict corresponding
shifts in $m_e$ that are testable against independent Penning-trap
measurements, providing a precision test of the framework at
$10^{-9}$ and beyond.

In the Standard Model, $\aem$ and $m_e$ are independent inputs;
no such tracking relation exists.
The framework's prediction~\eqref{P4:eq:consistency_condition}
is therefore qualitatively new and falsifiable.
\end{remark}

\begin{remark}[Why WZW bosonization gives the $\overline{\mathrm{MS}}$ mass]
\label{P4:rem:msbar_origin}
Witten non-abelian bosonization identifies the WZW fermion mass with the
$\overline{\mathrm{MS}}$ mass at the bosonization scale, because the
bosonization dictionary is derived from the path integral with dimensional
regularisation and $\overline{\mathrm{MS}}$ subtraction~\cite{Witten1984}.
The pole mass differs by the on-shell self-energy
$\delta m/m = \aem/\pi$ at one loop, independent of any cutoff.
The $+0.22\%$ residual of the three-term spoke was therefore
not a geometric failure but a renormalisation-scheme effect:
the framework predicted the correct $\overline{\mathrm{MS}}$ mass,
and the remaining gap was the scheme conversion.
\end{remark}

%──────────────────────────────────────────────────────────────────────────────
\subsection{Summary of section results}
\label{P4:sec:t_matrix_summary}
%──────────────────────────────────────────────────────────────────────────────

The argument proceeds in five steps:

\begin{enumerate}[noitemsep]

\item
  \emph{The $0.71\%$ gap is not Virasoro characters, WZ topology, duality, or framing.}
  Each candidate is rigorously ruled out (Section~\ref{P4:sec:ruled_out}).

\item
  \emph{The spoke exponent decomposes as (conformal weight) $-$ (Casimir).}
  The matching condition fixes $t_{\rm open}=80/13$, and $4\pi = 64\pi/13 - 12\pi/13$
  exactly (Proposition~\ref{P4:prop:exponent_decomp}).

\item
  \emph{The $j=1/2$ spoke carrier accumulates a $T$-matrix phase $T_{1/2}/Q = 3/400$.}
  As it traverses the $Q=10$ topological sectors, the modular monodromy
  contributes $-T_{1/2}/Q$ to the spoke exponent (Theorem~\ref{P4:thm:t_matrix_spoke}).

\item
  \emph{The correction is an exact WZW invariant.}
  $T_{1/2}/Q = c/(24Q) = 3/400$ follows from the $k=3$ identity
  $h_{1/2}=c/12$ (Lemma~\ref{P4:lem:three_forms}), itself equivalent to the
  self-duality $S_{1/2,1/2}=S_{0,0}$ underlying the whole framework.

\item
  \emph{The three-term $\aInv$ formula is now unconditional.}
  The corrected spoke determines $\mWZW/m_e$ from condensate geometry,
  removing the last free parameter (Theorem~\ref{P4:thm:alpha_unconditional}).

\end{enumerate}

The complete proof chain of Paper~IV, with all arrows proved:
\begin{equation}
  \underbrace{\bst = 1/\mWZW^2}_{\rm exact}
  \;\Rightarrow\;
  \underbrace{\mWZW^2/m_0^2 = \tfrac{8\vp}{15}\cdot\tfrac{2R_0^2}{2R_0^2+1}}_{\rm exact}
  \;\Rightarrow\;
  \underbrace{m_e/\Lcond = \vp^{20}e^{4\pi-3/400}}_{0.22\%}
  \;\Rightarrow\;
  \underbrace{\aInv = \tfrac{360}{\vp^2}-\tfrac{k}{2\pi}+\tfrac{1}{9\vp^6}}_{0.0004\%}.
  \label{P4:eq:chain_final}
\end{equation}

\begin{corollary}[Unconditional four-term $\aInv$ formula]
\label{P4:cor:four_term_unconditional}
Since Theorem~\ref{P4:thm:alpha_unconditional} makes the three-term formula
unconditional, the four-term formula of Theorem~\ref{P4:thm:four_term_alpha}
is also unconditional:
\begin{equation}
  \boxed{
  \aInv
  = \frac{360}{\vp^2} - \frac{k}{2\pi}
  + \frac{1}{9\vp^6} - \frac{1}{36\pi\vp^6}
  = 137.035\,998\,486\ldots
  }
  \label{P4:eq:four_term_unconditional}
\end{equation}
with residual $6.9\times10^{-7}$ ($5\times10^{-7}\%$) from the PDG value
$137.035\,999\,177$, a factor-of-$700$ improvement over the three-term result.
\end{corollary}

\begin{proof}
Theorem~\ref{P4:thm:four_term_alpha} derives the fourth term
$-1/(36\pi\vp^6)$ as a Chern--Simons correction to the one-loop QED
screening of Term~3.
That derivation builds on the three-term formula but does not invoke
the conditionality of $\mWZW/m_e$: the CS correction is computed from
the ratio $1/(4\pi)$ of the CS coupling to the WZW level, which is
an exact algebraic datum of the condensate.
Since Theorem~\ref{P4:thm:alpha_unconditional} establishes the three-term
formula unconditionally, its four-term extension
(Theorem~\ref{P4:thm:four_term_alpha}) inherits that status.
\qed
\end{proof}

%══════════════════════════════════════════════════════════════════════════════
\section{The running of $\alpha$ and the $Z$-boson scale}
\label{P4:sec:alpha_running}
%══════════════════════════════════════════════════════════════════════════════

The three-term formula predicts $\aInv = 137.035\,998$ at the \emph{electron
mass scale} $m_e$.  This is precisely where $\alpha$ is measured in
low-energy experiments: the electron anomalous magnetic moment ($g-2$),
atomic physics (hydrogen Lamb shift), and Thomson scattering.  A natural
question is how this prediction connects to the measured value
$\aInv(M_Z) \approx 128.9$ at the $Z$-boson scale.

\subsection{Two running mechanisms with opposite signs}

The framework contains two distinct running mechanisms that operate
simultaneously and in opposite directions, both acting between
$\mWZW$ and $m_e$. Neither runs from the geometric UV value
$\aInv_{\rm geo}=\vp^6/\sin^4(\pi/5)=150.3$: the step from $150.3$ to
the golden angle $360/\vp^2=137.508$ is $\Delta_1$, the exact,
computable difference between two independently-exact quantities
established in Result~3.1 of Paper~III~\cite{Paper3}%\ref{P3:prop:cascade}
, not a running effect (its deeper origin is Open Problem~4
of Paper~III~\cite{Paper3}%\ref{P3:sec:open}
). The two running mechanisms below act only on the
smaller, genuinely dynamical step from $137.508$ down to the measured
$137.036$.

\paragraph{WZW anti-screening (non-Abelian, dominant, $\mu > m_e$ regime).}
The $\mathrm{SU}(2)_3$ WZW condensate is non-Abelian.
Like QCD, it is asymptotically free: the effective coupling \emph{weakens}
at high energies and \emph{strengthens} at low energies.
The exact WZW beta function~\cite{KZ1984} gives:
\begin{equation}
  \frac{d\,\aInv}{d\log\mu} = -\frac{k}{2\pi}
  \label{P4:eq:wzw_beta_running}
\end{equation}
with negative sign (anti-screening: $\aInv$ \emph{decreases} at lower $\mu$).
This runs $\aInv$ from the golden-angle attractor $360/\vp^2=137.508$
down toward the measured value $137.036$, a shift of $-k/(2\pi)=-0.477$.
The WZW condensate acts on the electromagnetic coupling in the same way
that the gluon condensate acts on the strong coupling in QCD:
\emph{virtual flux-knot pairs in the condensate background enhance the
coupling at lower energies, reducing $\aInv$ from $137.508$ toward $137.036$.}
This is non-Abelian anti-screening, not the ordinary Abelian QED screening
from $e^+e^-$ pairs.

\paragraph{QED screening (Abelian, subleading, $m_e$ to $M_Z$ regime).}
Between the WZW decoupling scale $\mWZW$ and $m_e$, standard QED vacuum
polarization from virtual $e^+e^-$ pairs screens the bare charge.
This gives a positive (weakening) correction $\delta_{\rm QED} = +1/(9\vp^6)$
as the coupling runs from $\mWZW$ down to $m_e$.  This is the third term
in the formula.

\subsection{Boundary condition at $m_e$ and RGE above}

The framework thus provides a \emph{boundary condition}:
$\aInv(m_e) = 137.035\,998$.
Above $m_e$, the condensate decouples (Vainshtein mechanism: at densities
$\rho \gg 1/\bst$, the density-feedback suppression factor
$1/(1+\bst\rho)\to 0$ and the condensate decouples from local physics),
and the coupling evolves by standard QED and electroweak RGE.
The one-loop running from $m_e$ to $M_Z$ receives contributions from
the leptons ($e, \mu, \tau$), the five light quarks, and hadronic effects:
\begin{equation}
  \aInv(M_Z)
  \approx \aInv(m_e)
         - \frac{2}{3\pi}\sum_f N_c^f Q_f^2 \ln\frac{M_Z}{m_f}
         - \Delta\alpha^{-1}_{\rm had}
  \approx 128.9,
  \label{P4:eq:running_to_MZ}
\end{equation}
where $\Delta\alpha^{-1}_{\rm had} > 0$ is the hadronic vacuum polarization
contribution~\cite{PDG2024}, in full agreement with the measured
$\aInv(M_Z) = 128.9 \pm 0.1$.
 
The framework is therefore \emph{not} in competition with standard QED RGE
for the range $[m_e, M_Z]$.  It specifies the initial condition at $m_e$
with $5\times10^{-7}\%$ accuracy; the well-tested Standard Model RGE then carries
the coupling up to higher energies.

\subsection{The two-level picture}

\begin{center}
\renewcommand{\arraystretch}{1.3}
\begin{tabular}{lcll}
\toprule
Energy range & Running type & $\aInv$ direction & Carrier \\
\midrule
$\Lambda_{\rm cond} \to \mWZW$ & not running ($\Delta_1$, exact) & $150.3 \to 137.508$ (decreases) & --- \\
$\mWZW \to m_e$ & WZW anti-screening & $137.508 \to 137.036$ (decreases) & Virtual flux-knot pairs \\
$m_e \to M_Z$ & QED screening & $137.036 \to 128.9$ (decreases) & Virtual $e^+e^-$, $\mu^+\mu^-$, hadrons \\
\bottomrule
\end{tabular}
\end{center}

\noindent
The first row is not a running effect (Result~3.1 of Paper~III~\cite{Paper3}%\ref{P3:prop:cascade}
); it is included only to show where the geometric UV value sits
relative to the two genuine running steps below it.  Those two rows
show $\aInv$ decreasing, but for opposite physical reasons:
in the WZW regime, the coupling \emph{strengthens} at lower energies
(non-Abelian anti-screening increases $\alpha = 1/\aInv$); in the QED
regime, the coupling \emph{weakens} at higher energies (Abelian screening
decreases $\aInv$).  The two mechanisms are physically distinct and
operate in non-overlapping energy regimes once the condensate decouples.

\begin{remark}[The QED regime at low energies]
\label{P4:rem:qed_regime_clarification}
It may appear contradictory that QED screening increases $\aInv$
(weakens $\alpha$) at lower energies, while the standard textbook statement
is that ``virtual $e^+e^-$ pairs screen the charge, making $\alpha$
smaller at long distances.''  There is no contradiction: screening means
the \emph{apparent} charge is smaller when measured at lower energies
(larger distances), which means $\alpha$ is smaller, hence $\aInv$ is
\emph{larger} at lower energies.  Running \emph{down} from $\mWZW$ to $m_e$
therefore increases $\aInv$ by $+1/(9\vp^6) = +0.006$, as in the formula.
In contrast, the WZW anti-screening decreases $\aInv$ as we go to lower
energies (increases $\alpha$), the non-Abelian analogue of asymptotic
freedom.
\end{remark}

\newpage
%══════════════════════════════════════════════════════════════════════════════
\section{Status table}
\label{P4:sec:status}
%══════════════════════════════════════════════════════════════════════════════

\begin{center}
\begin{longtblr}[
  caption = {Consolidated Tier Table},
  label = {tab:tier},
] {
%{lp{3.5cm}cc}
  colspec = {X[2,p] X[1,p] X[1,c]  X[1.3,c]}, % Proportional widths
  rowhead = 1, % Repeat header on each page
  %row{1} = {font=\bfseries},
}
\toprule
Result & Status & Predicted & Deviation \\
\midrule
\SetCell[c=4]{l}{l}{\textit{Proved (this paper)}} \\
$V^*(R_0) = \vp$ exactly for all $R_0$ &
  \textcolor{darkgreen}{Theorem~\ref{P4:thm:verlinde_cancel}} & Exact & --- \\
$0.26\%$ Verlinde piece = grid noise ($O(h^2)$) &
  \textcolor{darkgreen}{Corollary~\ref{P4:cor:verlinde_exact}} & Exact & --- \\
$m_e/\TCMB$ is dimensionless geometric invariant &
  \textcolor{darkgreen}{Remark~\ref{P4:rem:single_input}} & Exact & --- \\
$m_e$ derivable from $\TCMB$ alone (up to $0.22\%$) &
  \textcolor{darkgreen}{Remark~\ref{P4:rem:single_input}} & $0.5121\,\mathrm{MeV}$, $O(R_0^{-2})$ & $0.22\%$ \\
$m_e$ derived (not assumed) from WZW spectrum &
  \textcolor{darkgreen}{Theorem~\ref{P4:thm:me_lightest}} & $0.5121\,\mathrm{MeV}$, $O(R_0^{-2})$ & $0.22\%$ \\
Full spoke-mass spectrum $m_j/\Lcond = \vp^{20}e^{4\pi-T_j/Q}$ &
  \textcolor{darkgreen}{Theorem~\ref{P4:thm:spoke_mass_spectrum}} & Exact & --- \\
$k=3$ uniquely selected by $T_{1/2}=c/24$ &
  \textcolor{darkgreen}{Theorem~\ref{P4:thm:casimir_selects}} & Exact & --- \\[4pt]
\SetCell[c=4]{l}{\textit{Proved (this paper)}} \\
$\bst = 1/\mWZW^2$ & \textcolor{darkgreen}{Theorem~\ref{P4:thm:compton}} & Exact & --- \\
$\sin^2\!\theta = 1/2$ (locked Hopf angle) & \textcolor{darkgreen}{Lemma~\ref{P4:lem:locked_angle}} & Exact & --- \\
$V(0;R_0)=\tfrac{15}{8}(1+\tfrac{1}{2R_0^2})$ & \textcolor{darkgreen}{Lemma~\ref{P4:lem:virial_factor}} & Exact & --- \\
$\mWZW^2/m_0^2 = (8\vp/15)\cdot 2R_0^2/(2R_0^2{+}1)$ & \textcolor{darkgreen}{Theorem~\ref{P4:thm:mass_renorm}} & Exact & --- \\
$k=3$ forces exactly 3 generations & \textcolor{darkgreen}{Theorem~\ref{P4:thm:three_gen}} & Exact & --- \\
$j_{1/2}\times j_{1/2}=j_0\oplus j_1$; $j_0$ ground state & \textcolor{darkgreen}{Theorem~\ref{P4:thm:spoke_ground}} & Exact & --- \\
$\vp^{2Q}$ factor in spoke: Verlinde + $Q=10$ & \textcolor{darkgreen}{Proved} & Exact & --- \\
Hub state $|B_\mathrm{hub}\rangle = |B_{3/2}\rangle$; $\Delta S_\mathrm{CS}=4\pi$ & \textcolor{darkgreen}{Theorem~\ref{P4:thm:second_spoke}} & Exact & --- \\
$m_e/\Lcond = \vp^{20}e^{4\pi}$ (leading order) & \textcolor{darkgreen}{Theorem~\ref{P4:thm:second_spoke}} & $0.97\%$ & --- \\
Self-duality $h_{1/2}=c/12$; $T_{1/2}=c/24=3/40$ & \textcolor{darkgreen}{Lemma~\ref{P4:lem:self_dual}} & Exact & --- \\
$T_{1/2}/Q = c/(24Q) = 3/400$ & \textcolor{darkgreen}{Lemma~\ref{P4:lem:three_forms}} & Exact & --- \\
$V^* = \vp$ & \textcolor{darkgreen}{Proved} & Exact & --- \\
$\Jfour/\Ja = 2^{4/3}/\vp^5$ & \textcolor{darkgreen}{Proved} & Exact & --- \\
$m_e/\Lcond = \vp^{20}e^{4\pi-3/400}$ (T-matrix corrected) & \textcolor{darkgreen}{Theorem~\ref{P4:thm:t_matrix_spoke}} & $0.5121\,\mathrm{MeV}$ & $0.22\%$ \\
$m_e/\Lcond = \vp^{20}e^{4\pi-3/400-\aem/\pi}$ (scheme corrected) & \textcolor{darkgreen}{Theorem~\ref{P4:thm:msbar_pole}} & $0.5110\,\mathrm{MeV}$ & $0.013\%$ \\
$\delta m_e/m_e = -\delta\aem/(\pi\aem)$ (tracking relation) & \textcolor{darkgreen}{Remark~\ref{P4:rem:alpha_me_tracking}} & Exact & --- \\
$\vEW = \Lcond\cdot\vp^{20}e^{49\pi/6}$ (leading) & \textcolor{darkgreen}{Corollary~\ref{P4:cor:vew}} & $249.76\,\mathrm{GeV}$ & $1.45\%$ \\
$\vEW = \Lcond\cdot\vp^{20}e^{49\pi/6-3/400}$ (corrected) & \textcolor{darkgreen}{Corollary~\ref{P4:cor:vew_corrected}} & $247.89\,\mathrm{GeV}$ & $0.69\%$ \\
\SetCell[c=4]{l}{\textit{Approximate identity (this paper)}} \\
$b^* \approx 3\vp^6/800$ & Proposition~\ref{P4:prop:bstar}, $0.21\%$ & $3\vp^6/800\approx0.0673$ & $0.21\%$ \\[4pt]
\SetCell[c=4]{l}{\textit{From \cite{Paper1,Paper2,Paper3}, recalled}} \\
$\aInv = 360/\vp^2 - k/(2\pi) = 137.030$ & \textcolor{darkgreen}{Proved in Paper~III~\cite{Paper3}} & $137.030$ & $0.004\%$ \\
$\aInv = 360/\vp^2 - k/(2\pi) + 1/(9\vp^6) = 137.036\,49$ & \textcolor{darkgreen}{Thm~\ref{P4:thm:three_term_alpha} (conditional in~\cite{Paper3}; unconditional: Thm~\ref{P4:thm:alpha_unconditional})} & $137.036\,49$ & $0.0004\%$ \\
$\aInv = 360/\vp^2 - k/(2\pi) + 1/(9\vp^6) - 1/(36\pi\vp^6) = 137.035\,998$ & \textcolor{darkgreen}{Thm~\ref{P4:thm:four_term_alpha}; unconditional: Cor.~\ref{P4:cor:four_term_unconditional}} & $137.035\,998$ & $5\times10^{-7}\%$ \\[4pt]
\bottomrule
\end{longtblr}
\end{center}

%══════════════════════════════════════════════════════════════════════════════
\section{Open problems}
\label{P4:sec:open}
%══════════════════════════════════════════════════════════════════════════════

\begin{enumerate}[noitemsep]

\item \textbf{[Resolved in Paper~V~\cite{Paper5}]}
The two-mechanism $\aInv$ pattern established here generalises to the
CP-violating phase: the dominant WZW spin sum overshoots at $76.5^\circ$
and the subleading $E_8$ McKay correction with the same condensate phase
quantum $1/9$ corrects to $\delta_{\rm CP} = 65.0^\circ$ ($0.12\sigma$
from PDG).  The weight $Q/(k\,h(E_8)) = 1/9$ is the same $1/(9\vp^6)$
denominator structure as the QED correction here.

\item \textbf{[Medium --- not needed for main results]} The transcendental correction to $b^* = 3\vp^6/800$ is established to all orders by the Newton formula~\eqref{P4:eq:bstar_newton}, which is exact.  The remaining task is to 
obtain a compact closed-form correction rather than the implicit transcendental expression. The error of $0.21\%$ in $b^*$ propagates to only $0.11\%$ in $m_\mathrm{WZW}$, so this is not blocking any main result.

\item \textbf{[Resolved in Paper~XII~\cite{Paper12}] --- path to single-input framework}
The profile normalisation conjecture
$\vp^9\sqrt{\Ja\Jfour} = Q = 10$ is proved in Paper~XII
via the two-route equivalence: the conjecture is equivalent to the
thermodynamic equilibrium condition
$E_{\rm fb}/V_{\rm torus} = \rho_{\rm CMB} = 10\,\Lcond^4$
(Theorems~A and E of Paper~XII~\cite{Paper12}).
This resolves both consequences stated here:

\begin{itemize}[noitemsep]
  \item The framework is now a \emph{one-input} model:
    $m_e$ is derived from $\TCMB$ alone
    (chain $\TCMB\to\rho_{\rm CMB}=10\to\vp^9\sqrt{\Ja\Jfour}=10\to m_e$).
  \item The spoke formula, $J_4/J_{2a}=2^{4/3}/\vp^5$, and
    $\alpha^{-1}$ hold unconditionally.
\end{itemize}

\item \textbf{[Experimental]} \textit{Material realization of the \texorpdfstring{$Q=1$}{Q=1} Hopfion.}
The framework's $Q=1$ topological soliton (the electron in the cosmic context)
requires a material analog with Hopf linking number 1 in its nodal-line structure.
$\mathrm{Li_2NaN}$ has been identified theoretically as a Hopf-link nodal-line
semimetal~\cite{Li2NaN2017,Li2NaN2018}: its two nodal loops in the Brillouin zone
are mutually interlocked with linking number $\pm 1$, realizing $Q=1$ Hopf charge.
This is topologically distinct from ZrSiS (unlinked loops, $Q=0$), which provides
the framework's symmetry structure ($k=3$, $2I$, $Q=10$) but not the Hopf linking.

The natural material hierarchy is:
\begin{equation*}
  \underbrace{\mathrm{ZrSiS}}_{\substack{Q=0,\;\mathbb{Z}_2\\\text{symmetry provider}}}
  \;\longrightarrow\;
  \underbrace{\mathrm{Li_2NaN}}_{\substack{Q=1,\;\text{link}\,{=}\pm1\\\text{Q=1 soliton}}}
  \;\longrightarrow\;
  \underbrace{\text{cosmic condensate}}_{\substack{Q=2,10\\\text{this framework}}}
\end{equation*}
The Hopf-link crossing gap $\Delta_{\mathrm{Li_2NaN}}$ in Li$_2$NaN sets the
natural energy unit for $Q=1$ quanta in the material.  Its ratio to the ZrSiS
spin-orbit coupling gap $\Delta_\mathrm{SOC}$ provides a dimensionless material
analog of $\mWZW/\Lcond$.  Precise ARPES measurement of $\Delta_{\mathrm{Li_2NaN}}$
(not yet reported as of 2026~\cite{Neupane2016,Lodge2024}) would constitute a
falsifiable experimental test of the framework's topological structure independent
of and complementary to the $\alpha$ calculation.

\end{enumerate}

%══════════════════════════════════════════════════════════════════════════════
\section{Summary}
\label{P4:sec:summary}
%══════════════════════════════════════════════════════════════════════════════

We have proved that the density-feedback Stieltjes parameter equals
the squared inverse WZW fermion mass (Theorem~\ref{P4:thm:compton}).
Two exact geometric lemmas underpin the mass renormalization:
Lemma~\ref{P4:lem:locked_angle} shows that the BS profile gradient makes
an exact $\pi/4$ angle with the radial direction at every cross-sectional
point (a topological invariant of the Hopf fibration), implying the
exact identity $\sin^2\!f/\rho^2 = \kappa_\mathrm{tube}/2$.
Lemma~\ref{P4:lem:virial_factor} gives the closed-form
$V(0;\,R_0) = \tfrac{15}{8}(1+\tfrac{1}{2R_0^2})$.
From these, Theorem~\ref{P4:thm:mass_renorm} gives the exact 3D fermion mass
renormalization factor $\sqrt{(8\vp/15)\cdot 2R_0^2/(2R_0^2+1)}$.

The three-term and four-term formulas $\aInv$
are now \emph{unconditional} (Theorem~\ref{P4:thm:alpha_unconditional} and Corollary~\ref{P4:cor:four_term_unconditional}).
Its conditionality in an earlier version rested on a $0.71\%$ gap in the spoke
formula $m_e/\Lcond = \vp^{20}e^{4\pi}$ whose origin was unidentified.
Section~\ref{P4:sec:t_matrix_correction} closes this gap via a single exact WZW
datum: the modular $T$-matrix eigenvalue $T_{1/2} = c/24 = 3/40$ of the
$j=1/2$ spoke carrier.  As the carrier traverses the $Q=10$ topological
sectors, it accumulates a phase $T_{1/2}/Q = 3/400$, which corrects the spoke
exponent to $4\pi - 3/400$.  Four mechanisms that might have explained the gap
(Virasoro character corrections, non-Abelian WZ topology, open-closed duality,
and framing anomaly) are rigorously ruled out.  The identity $T_{1/2}=c/24$
is exact at $k=3$ by the self-duality $S_{1/2,1/2}=S_{0,0}$, the same algebraic
specialness that selects $k=3$, $Q=10$, and $\dim_q(1/2)=\vp$.

The corrected spoke $m_e/\Lcond = \vp^{20}e^{4\pi-3/400}=4.305\times10^9$
achieves $0.22\%$ accuracy versus the measured $4.296\times10^9$ at
grid spacing $h=0.12$.
Theorem~\ref{P4:thm:verlinde_cancel} (Section~\ref{P4:sec:verlinde_cancel})
now establishes that the $0.26\%$ Verlinde correction previously attributed
to $V^*(R_0=3)=1.61824\neq\vp$ is algebraically zero: the two
$R_0$-dependent factors in $V(0;\,R_0)$ and $V^*/V(0)$ are exact
reciprocals and cancel for all $R_0$.
The $+0.22\%$ at $R_0=3$ is identified as the one-loop QED
pole-mass to $\overline{\mathrm{MS}}$ scheme conversion
$\aem/\pi = 0.2323\%$ (Theorem~\ref{P4:thm:msbar_pole}):
the WZW bosonization computes $m_e$ in $\overline{\mathrm{MS}}$
while the measured $0.511\,\mathrm{MeV}$ is the pole mass.
Incorporating this, the corrected formula
$m_e = \Lcond\cdot\vp^{20}\cdot e^{4\pi-3/400-\aem/\pi}$
achieves $0.013\%$ accuracy.
Any residual is $O(R_0^{-2})$ from the 2D toroidal geometry;
the theoretical thin-torus residual is $0.04\%$.

As a consequence, the ratio $m_e/\TCMB$ is a dimensionless geometric
invariant of the $Q=2$ icosahedral condensate:
\[
  \frac{m_e}{\TCMB}
  = \left(\frac{\pi^2}{150}\right)^{\!\!1/4}
    \cdot\vp^{20}\cdot
    \exp\!\left(\frac{1600\pi-3}{400}\right)
  \approx 2.18\times10^9
  \quad(k_B = 1).
\]
The framework takes a single experimental input $\TCMB$; Paper~XII proves
the one-parameter chain without any second input.
The ratio $m_e/\TCMB$ is a dimensionless geometric invariant of the
condensate; as $\aem$ is refined, $m_e$ tracks via
$\delta m_e/m_e = -\delta\aem/(\pi\aem)\approx-43.6\,\delta\aem/\aem$
(Remark~\ref{P4:rem:alpha_me_tracking}).
The ratio $\mWZW/m_e = e^{\pi/(k\vp^6)}$ is derived from condensate
geometry alone at $0.22\%$ accuracy, making $\delta_{\rm QED}=1/(9\vp^6)$
derivable without any external input and the $\aInv$ formula genuinely unconditional.

The four-term formula has a transparent physical structure: the golden angle
$360/\vp^2$ is the geometric UV attractor of the icosahedral condensate; the
WZW anti-screening correction $-k/(2\pi)$ reflects the non-Abelian nature of
the condensate (virtual flux-knot pairs enhance the coupling at lower
energies, the analogue of asymptotic freedom in QCD); the one-loop QED
screening correction $+1/(9\vp^6)$ is the Abelian contribution from virtual
$e^+e^-$ pairs running from $\mWZW$ down to $m_e$; and the Chern--Simons
correction $-1/(36\pi\vp^6) = -T_3/(4\pi)$ is the reduction of the QED
vacuum polarisation by the condensate's non-Abelian magnetic flux
(Theorem~\ref{P4:thm:four_term_alpha}).  The four-term formula
predicts $\aInv = 137.035\,998\,486$, matching the PDG value $137.035\,999\,177$
to $5\times10^{-7}\%$.  Above $m_e$, standard QED and electroweak running
takes over, reaching $\aInv\approx128.9$ at $M_Z$ in full agreement with
experiment (Section~\ref{P4:sec:alpha_running}).

The bicycle-wheel analogy identifies the ratio $\vEW/\Lcond$ as
a topological spoke of the Hopf fibration.  The $\mathrm{SU}(2)_3$
representation theory fixes exactly three spokes (three generations;
Theorem~\ref{P4:thm:three_gen}), and the fusion rule
$j_{1/2}\times j_{1/2}=j_0\oplus j_1$ identifies the $Q=2$ Hopfion
condensate as the ground state of the spoke system (Theorem~\ref{P4:thm:spoke_ground}).
The hub boundary state is $|B_{3/2}\rangle$; the Chern-Simons action difference
$\Delta S_\mathrm{CS} = 4\pi$ gives the spoke exponent, corrected to
$4\pi-3/400$ by the $T$-matrix mechanism.
The full spoke-mass spectrum $m_j/\Lcond = \vp^{20}e^{4\pi-T_j/Q}$ is
now derived for all three primaries (Theorem~\ref{P4:thm:spoke_mass_spectrum}),
and the electron is identified as the $j=\tfrac{1}{2}$ spoke via the Witten
bosonization dictionary and the $k=3$ self-duality $T_{1/2}=c/24$
(Theorems~\ref{P4:thm:me_lightest} and~\ref{P4:thm:casimir_selects})---the first
derivation of $m_e$ from the WZW spectrum without any assumption on its value.
The electroweak scale is thereby
expressed as $\vEW = \Lcond\cdot\vp^{20}e^{4\pi-3/400-\aem/\pi}$ (accurate
to $0.013\%$, Corollary~\ref{P4:cor:vew_corrected}), with no Standard Model
parameters as input.

The condensed-matter underpinning of the framework predicts a material
hierarchy: ZrSiS ($Q=0$, symmetry provider) and $\mathrm{Li_2NaN}$
($Q=1$, Hopf-link semimetal) realize the topological precursors of the
cosmic $Q=1$ and $Q=2$ solitons.  Measurement of the Hopf-link crossing
gap in $\mathrm{Li_2NaN}$ constitutes a directly falsifiable prediction
of the framework (Open Problem~3).

The three-correction structure of $\aInv$ established here --- dominant
non-Abelian WZW anti-screening, subleading Abelian QED screening
with weight $1/(9\vp^6)$, and the Chern--Simons correction
$-1/(36\pi\vp^6)$ --- has a two-correction analogue in
the companion Paper~V~\cite{Paper5}, where the
CP-violating phase $\delta_{\rm CP}$ is derived by the same pattern:
dominant untwisted-sector topological spin overshoots at $76.5^\circ$,
and the subleading $E_8$~McKay twisted-sector correction with the same
condensate phase quantum $Q/(k\,h(E_8)) = 1/9$ brings the prediction to
$65.0^\circ$, within $0.12\sigma$ of the PDG value.

%─── References ───────────────────────────────────────────────────────────────
\begin{thebibliography}{99}

\bibitem{Paper1}
F.~Manfredi,
\textit{The Density-Feedback Faddeev--Niemi Hopfion:
  Exact Algebraic Predictions for Standard-Model Parameters},
Zenodo (2026).
\href{https://doi.org/10.5281/zenodo.19342027}{doi:10.5281/zenodo.19342027}

\bibitem{Paper2}
F.~Manfredi,
\textit{BPS Structure and Fixed-Point Theorem for the
  Density-Feedback Faddeev--Niemi Hopfion},
Zenodo (2026).
\href{https://doi.org/10.5281/zenodo.19363491}{doi:10.5281/zenodo.19363491}

\bibitem{Paper3}
F.~Manfredi,
\textit{Two Constants from One Knot: The Fine Structure Constant and
  the Linking Scale as Exact Predictions of the Icosahedral WZW Condensate},
Zenodo (2026).
\href{https://doi.org/10.5281/zenodo.19478629}{doi:10.5281/zenodo.19478629}

\bibitem{Paper5}
F.~Manfredi,
\textit{Colour Confinement, the QCD Scale, and Hadronic Structure
  from the Density-Feedback Hopfion},
Preprint (2026).
\href{https://doi.org/10.5281/zenodo.19638857}{doi:10.5281/zenodo.19638857}

\bibitem{Paper12}
F.~Manfredi,
\textit{Paper~XII: The Profile Normalisation Conjecture --- Proof via CMB Energy Identity and Two-Route Equivalence},
\href{https://doi.org/10.5281/zenodo.19646945}{doi:10.5281/zenodo.19646945}

\bibitem{Faddeev1997}
L.~D. Faddeev and A.~J. Niemi,
\textit{Stable knot-like structures in classical field theory},
Nature \textbf{387}, 58 (1997).
\href{https://doi.org/10.1038/387058a0}{doi:10.1038/387058a0}

\bibitem{Witten1984}
E.~Witten,
\textit{Non-abelian bosonization in two dimensions},
Commun.\ Math.\ Phys.\ \textbf{92}, 455 (1984).
\href{https://doi.org/10.1007/BF01215276}{doi:10.1007/BF01215276}

\bibitem{Verlinde1988}
E.~Verlinde,
\textit{Fusion rules and modular transformations in 2D conformal field theory},
Nucl.\ Phys.\ B \textbf{300}, 360--376 (1988).
\href{https://doi.org/10.1016/0550-3213(88)90603-7}{doi:10.1016/0550-3213(88)90603-7}

\bibitem{Hardy1949}
G.~H.~Hardy,
\textit{Divergent Series},
Oxford University Press, Oxford (1949); repr.\ AMS Chelsea Publishing (1991);
ISBN 978-1-4704-7785-1.
% Standard reference for Stieltjes asymptotic series and Borel summability.

\bibitem{PDG2024}
R.~L. Workman et al.\ (Particle Data Group),
\textit{Review of Particle Physics},
Prog.\ Theor.\ Exp.\ Phys.\ \textbf{2022}, 083C01 (2022);
updated 2024.
\href{https://doi.org/10.1093/ptep/ptac097}{doi:10.1093/ptep/ptac097}

\bibitem{KZ1984}
V.~G. Knizhnik and A.~B. Zamolodchikov,
\textit{Current algebra and Wess--Zumino model in two dimensions},
Nucl.\ Phys.\ B \textbf{247}, 83 (1984).
\href{https://doi.org/10.1016/0550-3213(84)90374-2}{doi:10.1016/0550-3213(84)90374-2}

\bibitem{Cardy1989}
J.~L. Cardy,
\textit{Boundary conditions, fusion rules and the Verlinde formula},
Nucl.\ Phys.\ B \textbf{324}, 581 (1989).
\href{https://doi.org/10.1016/0550-3213(89)90521-X}{doi:10.1016/0550-3213(89)90521-X}

\bibitem{Ishibashi1989}
N.~Ishibashi,
\textit{The boundary and crosscap states in conformal field theories},
Mod.\ Phys.\ Lett.\ A \textbf{4}, 251 (1989).
\href{https://doi.org/10.1142/S0217732389000320}{doi:10.1142/S0217732389000320}

\bibitem{Witten1989}
E.~Witten,
\textit{Quantum field theory and the Jones polynomial},
Commun.\ Math.\ Phys.\ \textbf{121}, 351 (1989).
\href{https://doi.org/10.1007/BF01217730}{doi:10.1007/BF01217730}

\bibitem{Neupane2016}
M.~Neupane, I.~Belopolski, M.~M. Hosen et al.,
\textit{Observation of topological nodal fermion semimetal phase in ZrSiS},
Phys.\ Rev.\ B \textbf{93}, 201104(R) (2016),
\href{https://doi.org/10.1103/PhysRevB.93.201104}{doi:10.1103/PhysRevB.93.201104}

\bibitem{Lodge2024}
M.~S. Lodge, E.~Marcellina, Z.~Zhu, X.-P. Li, D.~Kaczorowski,
M.~S. Fuhrer, S.~A. Yang, B.~Weber,
\textit{Symmetry-selective quasiparticle scattering and electric field
tunability of the ZrSiS surface electronic structure},
Nanotechnology \textbf{35}, 195704 (2024),
\href{https://doi.org/10.1088/1361-6528/ad2639}{doi:10.1088/1361-6528/ad2639}

\bibitem{Li2NaN2017}
Z.~Yan, R.~Bi, H.~Shen, L.~Lu, S.-C. Zhang, Z.~Wang,
\textit{Nodal-link semimetals},
Phys.\ Rev.\ B \textbf{96}, 041103(R) (2017).
\href{https://doi.org/10.1103/PhysRevB.96.041103}{doi:10.1103/PhysRevB.96.041103}

\bibitem{Li2NaN2018}
W.~Chen, H.-Z. Lu, J.-M. Hou,
\textit{Topological semimetals with a double-helix nodal link},
Phys.\ Rev.\ B \textbf{96}, 041102(R) (2017);
$\mathrm{Li_2NaN}$ identified as Hopf-link semimetal
with nodal linking number $\pm 1$.
\href{https://doi.org/10.1103/PhysRevB.96.041102}{doi:10.1103/PhysRevB.96.041102}

\bibitem{Hanneke2008}
D.~Hanneke, S.~Fogwell, and G.~Gabrielse,
\textit{New Measurement of the Electron Magnetic Moment and the Fine Structure Constant},
Phys.\ Rev.\ Lett.\ \textbf{100}, 120801 (2008).
\href{https://doi.org/10.1103/PhysRevLett.100.120801}{doi:10.1103/PhysRevLett.100.120801}

\end{thebibliography}

\end{document}